% !TeX root = ../Book2.tex
\BookPart{变分法、弱解与椭圆方程}{b2:part:03}
% 本编导读（章级标题，不计入编号）
\BookPartGuide

上一编的显式核和函数展开在规则区域上十分有效；当区域不规则、系数随位置变化，或源项只具有有限可积性时，直接公式往往无法取得。本编改从函数空间中的能量和测试等式出发，依次回答三个问题：解是否存在，是否唯一并稳定依赖数据，存在的弱解又有多光滑。

\ProseParagraphBreak

\chapref{b2:ch:14}从允许扰动推导 Euler--Lagrange 方程，再用直接方法构造凸能量的极小元。\chapref{b2:ch:15}把对称能量推广为一般双线性型，完整证明 Lax--Milgram 定理，并处理有界可测系数。\chapref{b2:ch:16}以差商和截断估计研究正则性，分别说明系数、边界、维数和源项空间的作用。下图的箭头表示本编实际使用的前置关系。
\[
\begin{tikzcd}[column sep=2.7em,row sep=1.5em]
\text{第一卷弱收敛与 Sobolev 理论}\arrow[d]\arrow[dr]\arrow[drr]&&\\
\text{第14章}\arrow[r]&\text{第15章}\arrow[r]&\text{第16章}
\end{tikzcd}
\]
本编反复使用第一卷第八至九章的 Hilbert 与弱拓扑工具、第二十二章的零边界 Sobolev 空间，以及第二十四至二十五章的嵌入与紧性。调用处均写明卷号和结果编号。第十五章的 Fredholm 选读小节另用第一卷附录 C；它不参与第十六章普通强制椭圆方程的基本正则性证明。

\ProseParagraphBreak

阅读时宜始终把边界类和测试空间写在方程旁边。同一内部表达式配上零迹、零均值或自由边界测试，会成为不同问题。还应区分三个层次：第一变分为零只是驻点条件，强制型给出能量解，正则性估计才允许逐点读取微分方程。前一层的证明不应预借后一层才获得的导数。

\ProseParagraphBreak

第十六章末先完整证明可测系数标量方程的局部有界性，再以明确的证明概要介绍 Harnack 比较。普通阅读可以先接受\cref{b2:thm:16-measurable-harnack}，继续其后的 Hölder 推论；需要完整证明时，先补\secref{b2:sec:1701}中二进立方体与分化的工具，再读\cref{b2:thm:C-john-nirenberg}的指数尾，最后读\secref{b2:sec:H201}的幂迭代、对数测试及\cref{b2:thm:H2-linear-harnack}。这条补读路线无需先完成整套奇异积分理论。该比较及其 Hölder 推论不充当前面高阶正则性证明的前提。掌握本编后，可以继续第四编的实变调和分析，学习极大函数、奇异积分和频率分解怎样为更精细的方程估计提供工具。

% !TeX root = ../../Book2.tex
\chapter{变分法}
\label{b2:ch:14}
\BookChapterNumber{b2:ch:14}
% 本章摘要（不计入编号）
\begin{chapterabstract}
本章用函数空间中的微分研究积分泛函，推导标量与向量值 Euler--Lagrange 方程，并从边界项辨认自然边界条件。Dirichlet 能量说明凸性怎样把驻点变成唯一极小点；波动作用量则说明驻点未必是极小点。最后把第一卷的直接方法应用于梯度幂能量，逐项验证强制性、弱闭约束与弱下半连续性，得到在一般有界开集上的存在性，并说明弱方程与经典正则性之间仍有距离。
\end{chapterabstract}

% 本章目标（不计入编号）
\begin{learninggoals}
\begin{itemize}
\item\label{b2:item:14-goal-derivative}区分方向变分、Gâteaux 导数与 Fréchet 导数，检查积分下求导所需的控制。
\item\label{b2:item:14-goal-boundary}从允许扰动出发推导内部方程及自然边界条件，核对外法向的符号。
\item\label{b2:item:14-goal-energy}用精确能量差或第二变分判断极小性，区别作用量与守恒能量。
\item\label{b2:item:14-goal-direct}在 Sobolev 边界类中实施直接方法，指出强制性、弱紧性和下半连续性各自的作用。
\end{itemize}
\end{learninggoals}

一条连接固定端点的曲线怎样使长度最小？一张固定在边框上的薄膜怎样达到平衡？这类问题比较的是整个函数，而不只是有限个参数。微积分仍然有效，但导数、紧性和边界条件都必须放到合适的函数空间中。为避免把光滑形式计算误当作存在性证明，本章先推导候选者必须满足的方程，再独立构造极小元。

\ProseParagraphBreak

本章使用第一卷第九章的弱收敛与凸泛函理论、第二十二章的 Sobolev 空间，以及第二十四章的 Poincaré 不等式；具体调用处均写明结论。关于变分法的系统叙述，可进一步阅读 Bernard Dacorogna 的《Direct Methods in the Calculus of Variations》\cite{Dacorogna2008Direct}及 Lawrence C. Evans 的《Partial Differential Equations》\cite{Evans2010Partial}。本章只使用下面完整核验的凸能量模型，不预借更一般的非凸积分泛函下半连续性理论。

% !TeX root = ../../Book2.tex
\section{分离变量法}
\label{b2:sec:1201}

有限杆的温度既由初温决定，也由端点如何交换热量决定。空间模态要先满足端点约束，其时间衰减再由相应特征值决定。本节把分离变量所得级数放进已经建立的能量空间，随后用镜像表示及具体边界数据检验它的含义。

\subsection{有界区间}
\label{b2:subsec:120101}

热传导中，温度梯度引起从高温向低温的通量。设细杆的热容量密度为 \(w(x)>0\)、传导系数为 \(p(x)>0\)，通量取 \(-p(x)u_x\)。无热源时，对任意子区间的热量平衡给出 \(wu_t=(pu_x)_x\)。若再有与温度成正比的损耗 \(q(x)u\)，则方程成为
\[
 wu_t=(pu_x)_x-qu,
 \quad u_t+Au=0.
\]
在有限区间及\cref{b2:def:10-sl-operator}的连续系数、分离边界假设下，空间算子就是\chapref{b2:ch:10}建立的 Sturm--Liouville 算子。本节先取非负能量情形；例如 \(q\geq0\)，以及 Robin 系数非负。

\ProseParagraphBreak

对初温 \(f\in H=L^2(I,w\,\mathrm{d}x)\)，\cref{b2:prop:10-heat-energy-limit}给出
\begin{equation}\label{b2:eq:12-interval-expansion}
 u(t)=\sum_ne^{-\lambda_nt}f_ne_n,
 \quad f_n=\langle f,e_n\rangle_H.
\end{equation}
它在零时以 \(H\) 范数取得初值，正时间属于 \(D(A)\)，并在 \(H\) 中满足方程；还在 \(L^2((0,T);V)\) 中具有有限能量。连续系数并不使空间解自动无限次可微：当 \(p\) 只有连续性时，正则性应以 \(pu_x\) 和 \(D(A)\) 表述。下面转到常系数杆，才进一步证明所有空间阶数的经典光滑性。

\subsection{Fourier 级数}
\label{b2:subsec:120102}

在 \(0<x<\ell\) 上固定 \(\kappa>0\)，考虑 \(u_t=\kappa u_{xx}\)。齐次 Dirichlet 条件给出
\begin{equation}\label{b2:eq:12-dirichlet-series}
 u_D(t,x)=\sum_{n=1}^\infty f_ne^{-\kappa(n\pi/\ell)^2t}
 \sqrt{\dfrac2\ell}\sin\dfrac{n\pi x}{\ell}.
\end{equation}
齐次 Neumann 条件给出
\begin{equation}\label{b2:eq:12-neumann-series}
 u_N(t,x)=\overline f+
 \sum_{n=1}^\infty g_ne^{-\kappa(n\pi/\ell)^2t}
 \sqrt{\dfrac2\ell}\cos\dfrac{n\pi x}{\ell},
 \quad\overline f=\dfrac1\ell\int_0^\ell f.
\end{equation}
其中 \(f_n,g_n\) 分别为正弦和非零余弦系数。常数项来自零模态；它不随时间衰减。

\Needspace{6\baselineskip}
\begin{proposition}[正时间的经典解]\label{b2:prop:12-interval-smooth}
设 \(\ell,\kappa>0\)，\(u_D,u_N\) 分别为零 Dirichlet 和零 Neumann 边界的热模态级数\eqref{b2:eq:12-dirichlet-series}、\eqref{b2:eq:12-neumann-series}，系数取初值的标准正交正弦、余弦系数。若 \(f\in L^2((0,\ell);\mathbb C)\)，上面两个级数及其任意有限阶 \(x,t\) 导数在 \([\delta,T]\times[0,\ell]\) 上绝对一致收敛，其中 \(0<\delta<T\)。因此它们在正时间光滑，满足各自方程和边界条件。
\end{proposition}
\begin{proof}
求 \(j\) 次空间导数及 \(m\) 次时间导数只会把第 \(n\) 项乘以至多 \(C_{j,m}n^{j+2m}\)。Cauchy--Schwarz 给出统一可和上界
\[
 \sum_{n\geq1}|f_n|n^{j+2m}e^{-c\delta n^2}
 \leq\left(\sum_{n\geq1}|f_n|^2\right)^{1/2}
 \left(\sum_{n\geq1}n^{2j+4m}e^{-2c\delta n^2}\right)^{1/2}<\infty.
\]
余弦系数的证明相同。每个阶数都有相应一致收敛，故逐项微分合法。有限项已经满足方程和边界条件，极限保持它们。
\end{proof}

例如两端温度固定为零，初温在杆内恒为 \(1\)，则
\[
 u_D(t,x)=\dfrac4\pi\sum_{j=0}^{\infty}
 \dfrac{e^{-\kappa((2j+1)\pi/\ell)^2t}}{2j+1}
 \sin\dfrac{(2j+1)\pi x}{\ell}.
\]
正时间的端点值为零，而零时的 \(L^2\) 初温仍是常数一。它不可能在闭区间上随着 \(t\downarrow0\) 一致趋于一，因为端点值始终为零；初边值在角点不相容并不妨碍平方能量解存在。

\subsection{边界条件}
\label{b2:subsec:120103}

周期热核提供另一种表示。沿用第二编的实线核 \(H_s(x)=(4\pi s)^{-1/2}e^{-x^2/(4s)}\)，定义
\begin{align*}
 K_D(t,x,y)&=\sum_{m\in\mathbb Z}
 \left(H_{\kappa t}(x-y+2m\ell)-H_{\kappa t}(x+y+2m\ell)\right),\\
 K_N(t,x,y)&=\sum_{m\in\mathbb Z}
 \left(H_{\kappa t}(x-y+2m\ell)+H_{\kappa t}(x+y+2m\ell)\right).
\end{align*}
则 \(u_D(t,x)=\int_0^\ell K_D(t,x,y)\cdot f(y)\,\mathrm{d}y\)，Neumann 情形用 \(K_N\)。两式分别来自 \(f\) 的奇延拓、偶延拓及 \(2\ell\) 周期化；将卷积在每个周期分割后换元即得。高斯尾保证对正时间紧区间的一切求和与微分可交换。第二编的周期 Fourier 系数唯一性证明它们与式\eqref{b2:eq:12-dirichlet-series}、式\eqref{b2:eq:12-neumann-series}一致。

\ProseParagraphBreak

这两个空间表达式在小时间有直接用途：附近的镜像点控制边界影响，远处项指数小。大时间则应使用特征级数，只保留低频项。由于奇延拓与偶延拓在有限 \(p\) 范数中有界，周期热核逼近还给出 \(f\in L^p(0,\ell)\)、\(1\leq p<\infty\) 时的初值收敛。Dirichlet 的闭区间一致收敛要求 \(f\) 连续且两端为零；Neumann 情形对任意连续 \(f\) 都成立，零时不要求导数边界相容。

\ProseParagraphBreak

非齐次端点先用提升处理。例如 \(u(t,0)=a(t)\)、\(u(t,\ell)=b(t)\)，其中 \(a,b\in C^1\)，令
\[
 r(t,x)=a(t)\left(1-\dfrac x\ell\right)+b(t)\dfrac x\ell,
 \quad v=u-r.
\]
那么 \(v\) 的边界为零，且 \(v_t-\kappa v_{xx}=-r_t\)。记源项正弦系数为 \(F_n(t)\)，逐模态公式为
\[
 v_n(t)=e^{-\kappa\lambda_nt}v_n(0)
 +\int_0^te^{-\kappa\lambda_n(t-s)}F_n(s)\,\mathrm{d}s.
\]
对 \(F\in C([0,T];L^2)\)，积分项的平方范数由 \(\int_0^t\|F(s)\|_2\,\mathrm{d}s\) 控制，有限截断因紧时间区间上的连续像为紧集而一致逼近。因而该公式给出 \(C([0,T];L^2)\) 中的解；不能只因源项连续就额外断言所有高阶导数连续。

\ProseParagraphBreak

非齐次端温还可给出一个完全显式的周期响应。在半直线 \(x>0\) 上取端温 \(u(t,0)=\cos(\omega t)\)，其中 \(\omega>0\)。令 \(\alpha=\sqrt{\omega/(2\kappa)}\)，则
\[
 u(t,x)=e^{-\alpha x}\cos(\omega t-\alpha x)
\]
满足热方程与这个边界条件。验证时先用复值函数 \(e^{\mathrm{i}\omega t-(1+\mathrm{i})\alpha x}\)：其时间导数乘以 \(\mathrm{i}\omega\)，二阶空间导数乘以 \((1+\mathrm{i})^2\alpha^2=\mathrm{i}\omega/\kappa\)，再取实部即可。它的初温是 \(e^{-\alpha x}\cos(\alpha x)\)，因此此式描述已经建立的周期响应；若改给另一份初温，还须加上零端温差解，不能直接沿用同一式子。距离增加既压低振幅，也使相位推迟。振幅降为端点的 \(e^{-1}\) 倍所需距离为 \(\alpha^{-1}=\sqrt{2\kappa/\omega}\)，所以快速变化的端温主要影响靠近端点的一段区域。

\ProseParagraphBreak

Robin 条件的解仍用式\eqref{b2:eq:12-interval-expansion}，但总热量还受内部损耗影响。为将边界通量写得有意义，取有限区间 \(I=(a,b)\)、实连续系数 \(p,w,q\in C([a,b])\)、\(p,w>0\)。在开时间区间 \(J\subset(0,\infty)\) 上，设
\[
 u\in C^1(J;L^2(I,w\,\mathrm{d}x))\cap C(J;H^1(I)),
 \qquad pu_x\in C(J;H^1(I)),
\]
并在 \(L^2(I)\) 中满足 \(wu_t=(pu_x)_x-qu\)。此时通量 \(pu_x\) 有绝对连续代表及两个端点的迹，而积分是 \(L^2(I,w\,\mathrm{d}x)\) 上的连续线性泛函。因此可以先对时间求导，再用空间基本定理，得到
\begin{equation}\label{b2:eq:12-weighted-heat-balance}
 \frac{\mathrm{d}}{\mathrm{d}t}\int_a^b w(x)u(t,x)\,\mathrm{d}x
 =[p(x)u_x(t,x)]_a^b-\int_a^b q(x)u(t,x)\,\mathrm{d}x,
 \qquad t\in J.
\end{equation}
前面的谱解在任意闭正时间带上具有这些强解性质：指数衰减控制 \(Au\) 及 \(u_t\) 的 \(H\) 范数连续性。由\cref{b2:lem:10-form-coercive}及 \((pv')'=qv-wAv\)，还有
\[
 \|v\|_{H^1(I)}+\|pv'\|_{H^1(I)}
 \leq C\bigl(\|v\|_H+\|Av\|_H\bigr),\qquad v\in D(A).
\]
其中第一项先由移位能量控制，通量的零阶范数再由 \(v'\) 控制，一阶范数由所示方程控制。因此 \(u\) 与通量也具有刚才要求的 \(H^1\) 时间连续性。这里不需要把仅连续的 \(p\) 作经典微分，也不要求一般初温在零时具有边界导数。

\ProseParagraphBreak

式\eqref{b2:eq:12-weighted-heat-balance}的第一项是边界净输入，第二项是内部吸收或产生。Neumann 条件只消除前一项；即使边界绝热，内部反应也会改变总热量。例如在 \((0,\ell)\) 上取 \(p=w=q=1\)、齐次 Neumann 条件及初温 \(f=1\)，则 \(u(t,x)=e^{-t}\)，加权总热量为 \(\ell e^{-t}\)，并不守恒。

对无内部反应的纯热方程，取 \(w=1\)、\(p=\kappa>0\)、\(q=0\)，才得到
\[
 \frac{\mathrm{d}}{\mathrm{d}t}\int_0^\ell u(t,x)\,\mathrm{d}x
 =\kappa\bigl(u_x(t,\ell)-u_x(t,0)\bigr).
\]
零 Neumann 条件使总热量守恒。若使用 \(-\kappa u_x(t,0)+h_0u(t,0)=0\)、\(\kappa u_x(t,\ell)+h_\ell u(t,\ell)=0\)，其中 \(h_0,h_\ell\geq0\)，右端为 \(-h_0u(t,0)-h_\ell u(t,\ell)\)；它对非负温度不大于零。Dirichlet 固定端温却不固定通量，不能由齐次边界推出总热量不变。

\subsection{长时间行为}
\label{b2:subsec:120104}

\Needspace{6\baselineskip}
\begin{proposition}[趋向零模态]\label{b2:prop:12-long-time}
设 \(I=(a,b)\) 为有限区间，\(p,w,q\in C([a,b];\mathbb R)\)、\(p,w>0\)，令 \(A\) 为\cref{b2:def:10-sl-operator}中每端取一个实分离 Dirichlet 或 Robin 条件的算子，其作用为 \(Au=w^{-1}(-(pu')'+qu)\)，\(H=L^2(I,w\,\mathrm dx;\mathbb C)\)，并假设 \(A\geq0\)。对 \(f\in H\)，令 \(u(t)=e^{-tA}f=\sum_ne^{-\lambda_nt}\langle f,e_n\rangle_He_n\)，其中 \((e_n,\lambda_n)\) 是按特征值递增排列的标准正交谱基。\(P_0\) 为到 \(\ker A\) 的正交投影；若没有零特征值则 \(P_0=0\)。记最小正特征值为 \(\lambda_+\)。则对 \(t\geq0\)，
\[
 \|u(t)-P_0f\|_H\leq e^{-\lambda_+t}\|f-P_0f\|_H.
\]
对给定初值 \(f\)，热解趋于零当且仅当 \(P_0f=0\)。所有 \(f\in H\) 的热解都趋于零，当且仅当 \(\ker A=\{0\}\)，这在当前非负正则有限区间条件下等价于 \(\lambda_1>0\)。特别地，对 \(I=(0,\ell)\)、\(A=-\kappa\partial_x^2\)、\(\ell,\kappa>0\)，零 Dirichlet 情形 \(P_0f=0\)，零 Neumann 情形 \(P_0f=\overline f=\ell^{-1}\int_0^\ell f\)；两者 \(\lambda_+=\kappa\pi^2/\ell^2\)。
\end{proposition}
\begin{proof}
将零模态从展开中分离，剩余每个系数的平方都乘以不大于 \(e^{-2\lambda_+t}\) 的因子，求和即得。最小正特征值存在由\cref{b2:thm:10-sl-spectral}保证。于是 \(u(t)\to P_0f\)，给定初值趋零的判据随之成立。若要求所有初值都趋零，取非零 \(f\in\ker A\) 就会得到不动解 \(u(t)=f\)，所以必须核为零；反向由已得估计成立。谱从最低特征值开始且趋于正无穷，非负性下核为零恰好意味着 \(\lambda_1>0\)。纯热 Dirichlet、Neumann 的结论由各自显式谱得到。
\end{proof}

对上述常系数 Dirichlet、Neumann 区间，这还给出 \(t\geq1\) 时的一致收敛：先演化至 \(t-1\)，再由固定时间一的核把 \(L^2\) 误差送到一致范数；Cauchy--Schwarz 与正时间可和尾给出固定的 \(L^2\to L^\infty\) 界。若初值恰与最低正模态正交，衰减速度由下一项非零系数决定。出现负特征值时，相应非零初始系数会产生指数增长，不能继续使用本命题。

\ProseParagraphBreak

对于一般连续系数，判断是否趋于零应检查整个零空间。一个容易核验的充分条件是：\(q\geq0\)、所有 Robin 系数 \(h_s\geq0\)，且至少有一个 Dirichlet 端点或一个严格正的 Robin 系数。事实上，若 \(v\in\ker A\)，则
\[
 0=\langle Av,v\rangle_H
  =\int_a^b\bigl(p|v'|^2+q|v|^2\bigr)\,\mathrm{d}x
    +\sum_{s\in R}h_s|v(s)|^2.
\]
各项非负，\(p>0\) 使 \(v'=0\) 几乎处处，故 \(v\) 为常数。Dirichlet 端点的零值或严格正的 Robin 边界项迫使这个常数为零。因此核为零，由命题得到所有初温的指数衰减。这个论证用到了 \(q\geq0\)，不能仅留下正 Robin 条件而删去对内部系数的控制。

下面的例子说明原因。取 \(I=(0,1)\)、\(p=w=1\)，令
\[
 v_0(x)=1+x(1-x)>0,\qquad
 q(x)=\frac{v_0''(x)}{v_0(x)}=-\frac2{1+x(1-x)},
\]
并在两端取严格正 Robin 系数 \(h_0=h_1=1\)，即
\[
 -v'(0)+v(0)=0,\qquad v'(1)+v(1)=0.
\]
函数 \(v_0\) 满足这两个条件，而且 \(Av_0=-v_0''+qv_0=0\)。还需核验这个例子没有违反非负性假设。对任意 \(v\in H^1(0,1;\mathbb C)\)，令 \(z=v/v_0\)。由于 \(v_0\) 光滑且有正下界，\(z\in H^1(0,1)\)，绝对连续函数的乘积法则给出
\[
 |v'|^2+\frac{v_0''}{v_0}|v|^2
 =v_0^2|z'|^2+\bigl(v_0v_0'|z|^2\bigr)'.
\]
对 \(x\) 积分，导数项成为
\[
 \left[\frac{v_0'}{v_0}|v|^2\right]_0^1
 =-|v(1)|^2-|v(0)|^2,
\]
恰好抵消能量形式中的两个正边界项。因此
\begin{equation}\label{b2:eq:12-robin-zero-mode-factorization}
 \begin{aligned}
 \mathfrak a(v,v)
 &=\int_0^1\bigl(|v'|^2+q|v|^2\bigr)\,\mathrm{d}x
       +|v(0)|^2+|v(1)|^2\\
 &=\int_0^1v_0^2\left|\left(\frac v{v_0}\right)'\right|^2\,\mathrm{d}x
 \geq0.
 \end{aligned}
\end{equation}
所以 \(A\geq0\)，却仍有非恒定零模态。取初温 \(f=v_0\)，就得到永不衰减的解 \(u(t,x)=v_0(x)\)。正 Robin 系数排除了非零常数零模态，但允许 \(q\) 有负值时，不能据此排除所有零模态。

\ProseParagraphBreak

最后回到 \((0,\ell)\) 上 \(u_t=\kappa u_{xx}\)、\(\kappa,\ell>0\) 的纯热模型，把边界的几种作用放在一起比较。周期边界把两个端点耦合为
\[
 u(t,0)=u(t,\ell),\qquad u_x(t,0)=u_x(t,\ell),
\]
因而不属于前面的分离边界条件，不能直接套用其特征值单重的编号。对 \(f\in L^2((0,\ell);\mathbb C)\)，将区间视为周长 \(\ell\) 的圆，独立使用圆周 Fourier 理论：令
\[
 e_k(x)=\ell^{-1/2}e^{2\pi\mathrm{i}kx/\ell},\qquad
 f_k=\int_0^\ell f(x)\overline{e_k(x)}\,\mathrm{d}x,
 \qquad k\in\mathbb Z.
\]
周期热解为
\begin{equation}\label{b2:eq:12-periodic-heat-series}
 u_{\mathrm{per}}(t,x)=\sum_{k\in\mathbb Z}
       f_ke^{-\kappa(2\pi k/\ell)^2t}e_k(x),\qquad t\geq0.
\end{equation}
Parseval 等式及控制收敛使它在零时以 \(L^2\) 范数取得初值；任意正时间带上的指数尾又允许所需的逐项微分。零频率给出 \(\overline f=\ell^{-1}\int_0^\ell f\)，其余频率至少为 \(2\pi/\ell\)，所以
\[
 \|u_{\mathrm{per}}(t)-\overline f\|_2
 \leq e^{-4\kappa\pi^2t/\ell^2}\|f-\overline f\|_2.
\]
周期边界的两个端点通量相消，因此总热量守恒。保正性也可独立看出：周期化 Gaussian
\[
 P_\ell(t,z)=\sum_{m\in\mathbb Z}
  \frac1{\sqrt{4\pi\kappa t}}
  e^{-(z+m\ell)^2/(4\kappa t)},\qquad t>0,
\]
严格为正，且 \(\int_0^\ell P_\ell(t,z)\,\mathrm{d}z=1\)。把 Fourier 系数的积分按周期换元，再用 Gaussian 的变换，得到其第 \(k\) 个系数为 \(\ell^{-1}e^{-\kappa(2\pi k/\ell)^2t}\)；这里系数采用 \(\ell^{-1}\int_0^\ell P_\ell(t,z)e^{-2\pi\mathrm{i}kz/\ell}\,\mathrm{d}z\) 的非单位约定。因此式\eqref{b2:eq:12-periodic-heat-series}也等于 \(\int_0^\ell P_\ell(t,x-y)f(y)\,\mathrm{d}y\)，非负初温产生非负解。

分离边界的纯热模型同样保正，但原因并非仅为谱非负。对实值解令 \(u_-=\max\{-u,0\}\)，在 Dirichlet 端点有 \(u_-=0\)，在 Robin 端点使用 \(-\kappa u_x(0)+h_0u(0)=\kappa u_x(\ell)+h_\ell u(\ell)=0\)，其中 \(h_0,h_\ell\geq0\)，Neumann 对应 \(h_s=0\)。负部属于 \(H^1\)，且其弱导数为 \(-u_x\mathbf1_{\{u<0\}}\)。对正时间强解，将方程与负部配对并分部积分，得到
\[
 \frac12\frac{\mathrm{d}}{\mathrm{d}t}\|u_-\|_2^2
 =-\kappa\|(u_-)_x\|_2^2
      -\sum_{s\in R}h_s|u_-(t,s)|^2\leq0.
\]
时间链式法则可由 \(s\mapsto s_-^2/2\) 的 Lipschitz 导数及 \(L^2\) 时间可微性验证；空间分部积分只用通量和负部的一阶弱导数。从 \(\varepsilon>0\) 积分到 \(t\)，再用 \(u(\varepsilon)\to f\) 及负部映射的 \(L^2\) 连续性令 \(\varepsilon\downarrow0\)，就得到非负初温始终非负。\secref{b2:sec:1203}将对负部方法及最大值原理作进一步说明。

\begin{table}[htbp]
\centering
\small
\setlength{\tabcolsep}{4pt}
\caption[纯热方程的边界、热量与长期平衡]{纯热方程的边界、热量与长期平衡。固定 \(\ell,\kappa>0\)，方程为 \(u_t=\kappa u_{xx}\)，初温 \(f\in L^2((0,\ell);\mathbb R)\)；总热量指 \(\int_0^\ell u\)，长期极限按 \(L^2\) 范数理解，\(\overline f=\ell^{-1}\int_0^\ell f\)。保正指 \(f\geq0\) 时 \(u(t)\geq0\)。Robin 采用正文的含 \(\kappa\) 通量约定。}
\label{b2:tab:12-boundary-heat-equilibrium}
\begin{tabularx}{\linewidth}{@{}>{\raggedright\arraybackslash}p{3.5cm}>{\raggedright\arraybackslash}X>{\raggedright\arraybackslash}p{1.1cm}>{\raggedright\arraybackslash}p{1.8cm}@{}}
\toprule
边界条件 & 总热量 & 保正 & 长期极限 \\
\midrule
两端零 Dirichlet & 一般不守恒 & 是 & \(0\) \\
两端零 Neumann & 守恒 & 是 & \(\overline f\) \\
Robin：\(h_0,h_\ell\geq0\)，至少一端严格正 & 非负初温时不增 & 是 & \(0\) \\
周期边界 & 守恒 & 是 & \(\overline f\) \\
\bottomrule
\end{tabularx}
\end{table}

\Cref{b2:tab:12-boundary-heat-equilibrium}只比较无内部反应的纯热方程。质量守恒、保正性与长期趋零是三个不同的问题：加入 \(q(x)u\) 后，须分别重新检查加权热量平衡、解的符号及 \(\ker A\) 和谱下界，不能仅凭某一个边界系数作判断。

% !TeX root = ../../Book2.tex
\section{Fourier 变换法}
\label{b2:sec:1202}

在全空间中，空间平移不会改变方程，离散特征级数因而换成连续频率积分。变换负责找出核，卷积与逼近恒等负责证明粗初值确实演化成解。外力也可逐时叠加，但积分公式本身并不自动赋予它任意阶的经典导数。

\subsection{热核}
\label{b2:subsec:120201}

全空间没有端点，也没有正的最低允许频率。Fourier 变换把每个连续频率单独演化。先对足够光滑且衰减的解形式计算，按照本卷的 \(2\pi\) 约定有
\[
 \partial_t\widehat u(t,\xi)
 =-4\pi^2\kappa|\xi|^2\widehat u(t,\xi),
 \quad\widehat u(t,\xi)=e^{-4\pi^2\kappa t|\xi|^2}\widehat f(\xi).
\]
\cref{b2:thm:03-gaussian}已经计算了这个乘子的逆变换。沿用第二编的热核记号，在 \(\mathbb R^d\) 上写作
\begin{equation}\label{b2:eq:12-gauss-heat}
 H_t(x)=(4\pi t)^{-d/2}e^{-|x|^2/(4t)},
 \quad t>0,
 \quad\widehat H_t(\xi)=e^{-4\pi^2t|\xi|^2}.
\end{equation}
方程扩散系数为 \(\kappa\) 时使用 \(H_{\kappa t}\)，不另换 Fourier 归一化。这个核非负、总积分为一；其空间宽度约为 \(\sqrt t\)，高度约为 \(t^{-d/2}\)。由直接微分有 \(\partial_tH_{\kappa t}=\kappa\Delta H_{\kappa t}\)。

\ProseParagraphBreak

这些计算可以独立验证卷积候选，不要求先假定未知解的变换存在。对于粗初值，求解的逻辑次序是先用核定义 \(u\)，再证明可微性和初值收敛；前面的频率微分只负责找到候选。第二编已经证明 Gauss 核的逼近恒等性质，下面把其参数取为 \(\sqrt{4\pi\kappa t}\)。

\subsection{卷积表示}
\label{b2:subsec:120202}

\Needspace{6\baselineskip}
\begin{theorem}[全空间热核解]\label{b2:thm:12-heat-convolution}
设整数 \(d\geq1\)、\(\kappa>0\)，\(H_s(x)=(4\pi s)^{-d/2}e^{-|x|^2/(4s)}\)（\(s>0\)、\(x\in\mathbb R^d\)）。设 \(1\leq p\leq\infty\)、\(f\in L^p(\mathbb R^d;\mathbb C)\)。则
\begin{equation}\label{b2:eq:12-heat-convolution}
 u(t,x)=\int_{\mathbb R^d}H_{\kappa t}(x-y)\cdot f(y)\,\mathrm{d}y
\end{equation}
对 \(t>0\)、\(x\in\mathbb R^d\) 有意义，且在 \((0,\infty)\times\mathbb R^d\) 上光滑，满足 \(u_t=\kappa\Delta u\)。当 \(t\downarrow0\) 且 \(p<\infty\) 时，\(u(t)\to f\) 于 \(L^p\)。若选定的代表 \(\widetilde f\) 有界一致连续，则 \(u(t)\to\widetilde f\) 在一致范数中成立；若选定 \(f\) 的一个有界代表 \(\widetilde f\)，且它在 \(x_0\in\mathbb R^d\) 连续，则 \(u(t,x_0)\to\widetilde f(x_0)\)。
\end{theorem}
\begin{proof}
热核及其每阶时空导数属于每个 \(L^r\)、\(1\leq r\leq\infty\)。固定 \(t\) 的正紧区间与 \(x\) 的紧集后，每个导数可由一个平移后仍快速衰减的函数统一控制。对 \(f\in L^p\)，Hölder 不等式给出绝对可积性；对差商使用同一控制，可以把每一阶导数移入积分。因此卷积光滑，且核的微分等式传给 \(u\)。这里在 \(p=1\) 用核导数的一致界，在 \(p=\infty\) 用其绝对可积性，两个端点都被覆盖。

有限 \(p\) 与一致连续代表的情形由\cref{b2:thm:03-approximate-identity}得到。卷积不因初值在零测集上的修改而改变。对选定代表的连续点，把 \(\widetilde f(x_0-y)-\widetilde f(x_0)\) 的积分分成 \(|y|<\delta\) 与补集。近处由连续性控制，远处由 \(2\sup_x|\widetilde f(x)|\) 乘热核尾质量控制，先选 \(\delta\) 再令 \(t\downarrow0\) 即得。这样取得的点值是所选连续代表的点值，而不是给 \(L^p\) 等价类额外指定一个任意值。
\end{proof}

初值也可以是一点质量，此时解是 \(H_{\kappa t}(x-y_0)\)，在 \(t\downarrow0\) 时只以测度或分布意义趋于 \(\delta_{y_0}\)。为了区别空间核与时空基本解，令 \(E(t,x)=H_{\kappa t}(x)\) 当 \(t>0\)，在 \(t<0\) 令 \(E=0\)。那么
\[
 (\partial_t-\kappa\Delta_x)E=\delta_{(0,0)}
\]
在时空分布中成立。确实，对紧支集光滑 \(\varphi(t,x)\)，从 \(t=\varepsilon\) 起分部积分，热方程抵消内部项，留下 \(\int H_{\kappa\varepsilon}(x)\cdot\varphi(\varepsilon,x)\,\mathrm{d}x\)。它由逼近恒等趋于 \(\varphi(0,0)\)。同时 \(E\) 在有限时间带的空间积分为该带长度，所以零时附近的局部可积性也已得到。

\ProseParagraphBreak

对非齐次方程，设整数 \(d\geq1\)、\(\kappa,T>0\)、\(1\leq p<\infty\)，\(f\in L^p(\mathbb R^d;\mathbb C)\)、\(F\in C([0,T];L^p(\mathbb R^d;\mathbb C))\)。由逐时叠加得到
\begin{equation}\label{b2:eq:12-duhamel}
 u(t)=H_{\kappa t}*f+
 \int_0^tH_{\kappa(t-s)}*F(s)\,\mathrm{d}s.
\end{equation}
在积分的右端按 \(H_0*F(t)=F(t)\) 理解。被积函数是连续的 \(L^p\) 值函数，积分可作为 Riemann 和的范数极限；Young 不等式给出
\[
 \left\|\int_0^tH_{\kappa(t-s)}*F(s)\,\mathrm{d}s\right\|_p
 \leq\int_0^t\|F(s)\|_p\,\mathrm{d}s.
\]
这保证源项趋于零初值。先对光滑紧支集的时空源项求导，再以连续 \(L^p\) 范数逼近一般 \(F\)，便在分布中得到 \(u_t-\kappa\Delta u=F\)。逼近可将时间分成有限小段，对有限个空间取值分别磨光，再以时间连续插值连接；上式使解同时收敛，测试函数的积分配对也连续。对于这里的一般连续 \(L^p\) 值源项，方程按分布意义成立；经典二阶空间正则性还需进一步的条件。

\subsection{高维热方程}
\label{b2:subsec:120203}

高维热核是各坐标一维核的乘积，但高维扩散的幅度不能从一维直观中省略维数。换元 \(x=\sqrt t\,z\) 给出
\[
 \|H_t\|_r=C_{d,r}t^{-\frac d2(1-1/r)},
 \quad1\leq r\leq\infty.
\]
例如有限质量初温满足 \(\|u(t)\|_\infty\leq(4\pi\kappa t)^{-d/2}\|f\|_1\)。没有任何质量损失也能使温度峰值降低，因为同一热量散布到越来越大的区域。

\ProseParagraphBreak

方向上的扩散速率也可以不同。若整数 \(d\geq1\)，\(B\) 是实对称正定的 \(d\times d\) 常矩阵，方程为 \(u_t=\sum_{i,j=1}^dB_{ij}\partial_i\partial_ju\)，则 \(t>0\) 时的相应核是
\[
 H_t^B(x)=\dfrac1{(4\pi t)^{d/2}\sqrt{\det B}}
 \exp\left(-\dfrac{x^{\mathsf T}B^{-1}x}{4t}\right).
\]
正交对角化 \(B\) 后，各方向分别使用一维扩散核，换元的行列式给出归一化因子。故核仍非负、质量为一，变换为 \(e^{-4\pi^2t\xi^{\mathsf T}B\xi}\)。正定性使每个方向都发生扩散。若 \(B\) 仅半正定且有零特征值，\(B^{-1}\) 与 \((\det B)^{-1/2}\) 就没有上述意义，不能在函数密度公式中直接把某个特征值置零。

\ProseParagraphBreak

例如在 \(\mathbb R^2\) 上取 \(B=\operatorname{diag}(1,0)\)，方程成为 \(u_t=u_{xx}\)，没有 \(y\) 方向的扩散。此时传播核是乘积测度
\[
 H_t^{(1)}(x)\otimes\delta_0(y),\qquad
 H_t^{(1)}(x)=(4\pi t)^{-1/2}e^{-x^2/(4t)},\quad t>0,
\]
它支撑在直线 \(y=0\) 上，不是相对于二维 Lebesgue 测度的可积函数密度。对 \(f\in L^2(\mathbb R^2;\mathbb C)\)，传播只沿 \(x\) 作卷积：
\[
 u(t,x,y)=\int_{\mathbb R}H_t^{(1)}(x-z)f(z,y)\,\mathrm{d}z.
\]
Fubini 定理使几乎每个 \(y\) 对应的初值切片属于 \(L^2(\mathbb R_x)\)，故这一维卷积有意义，并在分布意义下满足 \(u_t=u_{xx}\)。但它不提高横向正则性：取非负非零 \(\phi\in C_c^\infty(\mathbb R;\mathbb R)\)，令 \(f(x,y)=\phi(x)\mathbf1_{(0,1)}(y)\)，就有
\[
 u(t,x,y)=\bigl(H_t^{(1)}*\phi\bigr)(x)\mathbf1_{(0,1)}(y).
\]
前一个因子在 \(t>0\) 时光滑且处处为正，\(y=0,1\) 处的跳跃却仍保留。一般半正定矩阵经正交对角化后也如此：正特征值方向使用 Gaussian 核，零特征值方向保留点质量，不能从正方向的磨光推出全空间光滑化。

\ProseParagraphBreak

这一正定二次型公式可与 John K. Hunter、Bruno Nachtergaele 的《Applied Analysis》\cite[pp. 303--304]{Hunter2001Applied}对照，其变换采用不同的 \(2\pi\) 归一化。本卷已经在\secref{b2:sec:0502}换算过常数，这里直接使用一致的频率约定。

% !TeX root = ../../Book2.tex
\section{能量泛函}
\label{b2:sec:1403}

同一个偏微分方程可以通过能量的驻点条件出现，但“能量”一词并不总表示待极小化的非负量。椭圆方程常对应空间能量的极小值；波动方程对应时空作用量的驻点，同时另有随时间守恒的正能量。本节用这两个模型把几种不同作用区分开。

\subsection{Dirichlet 能量}
\label{b2:subsec:140301}

设 \(\Omega\subset\mathbb R^d\) 是有界开集，记 \(H_0^1(\Omega)=W_0^{1,2}(\Omega)\)。第一卷定义 22.2.1 将它定义为 \(C_c^\infty(\Omega)\) 在 \(H^1\) 中的闭包，这一定义不要求边界光滑。由第一卷定理 24.1.2 的零边界 Poincaré 不等式，\(\lVert v\rVert_2\le C_P\lVert\nabla v\rVert_2\)，所以 \(\lVert\nabla v\rVert_2\) 是这个空间上的等价 Hilbert 范数。

\ProseParagraphBreak

设 \(\Omega\subset\mathbb R^d\) 是有界开集，给定实值 \(f\in L^2(\Omega)\) 和实值 \(g\in H^1(\Omega)\)，把允许函数写为 \(g+H_0^1(\Omega)\)，考虑 \term[dirichlet-nengliang]{Dirichlet 能量}[Dirichlet energy]及线性外力项组成的泛函
\[
 E(u)=\frac12\int_\Omega|\nabla u|^2\,\mathrm dx-\int_\Omega fu\,\mathrm dx.
\]
仿射空间的写法精确表达“\(u\) 与 \(g\) 具有相同边界数据”：即使边界粗糙到尚未定义逐点边界迹，也可以使用它。

\Needspace{6\baselineskip}
\begin{proposition}[二次能量的驻点与极小点]
\label{b2:prop:14-dirichlet-minimum}
设整数 \(d\geq1\)，\(\Omega\subset\mathbb R^d\) 为非空有界开集，\(f\in L^2(\Omega;\mathbb R)\)、\(g\in H^1(\Omega;\mathbb R)\)，\(E(v)=\frac12\int_\Omega|\nabla v|^2-\int_\Omega fv\)。\(u\in g+H_0^1(\Omega;\mathbb R)\) 是 \(E\) 在该仿射空间上的极小点，当且仅当
\[
 \int_\Omega\nabla u\cdot\nabla h\,\mathrm dx=\int_\Omega fh\,\mathrm dx
 \qquad(h\in H_0^1(\Omega)).
\]
这样的极小点至多有一个，且对每个 \(h\in H_0^1(\Omega)\) 都有
\[E(u+h)-E(u)=\tfrac12\lVert\nabla h\rVert_2^2.\]
\end{proposition}
\begin{proof}
Cauchy--Schwarz 不等式保证所有积分有限。直接展开平方得
\[
 E(u+h)-E(u)=\int_\Omega\nabla u\cdot\nabla h-\int_\Omega fh
              +\tfrac12\lVert\nabla h\rVert_2^2.
\]
其中前两项构成连续线性泛函，最后一项是二阶余项，故一阶导数正是前两项。极小点的双侧变分为零，得到所列积分等式。反之，积分等式使能量差等于非负的梯度平方，所以 \(u\) 为全局极小点。若另一点 \(u+h\) 也极小，则 \(\nabla h=0\)；Poincaré 不等式推出 \(h=0\)。
\end{proof}

这一命题暂时没有断言极小点存在；它说明一旦构造出驻点，就无须再排查鞍点。二次展开还给出误差的精确解释：任何满足同一边界数据的近似函数，其能量超额恰为梯度误差平方的一半。与一般的第二变分局部判别相比，这是全局恒等式，原因在于能量严格二次且对允许方向一致凸。

\ProseParagraphBreak

在区间 \((0,1)\) 上取 \(f=1\)、\(g=0\)。函数 \(u(x)=x(1-x)/2\) 满足 \(-u''=1\) 和零端点条件；对光滑紧支集扰动分部积分，再用密度推广到 \(H_0^1\)，即可验证积分等式。因此它是唯一极小点，能量为 \(-1/24\)。负能量并不矛盾：梯度项非负，外力项却可以降低总泛函值；强制性要求的是大范数下的增长，而不是每一点的非负性。

\subsection{波动作用量与守恒能量}
\label{b2:subsec:140302}

设整数 \(d\geq1\)，\(\Omega\subset\mathbb R^d\) 为非空有界 \(C^1\) 开集，\(c,T>0\)，并考虑在 \([0,T]\times\overline\Omega\) 的开邻域内为 \(C^2\)、在侧边界取零值的实函数 \(u(t,x)\)。时空作用量为
\[
 \mathcal A(u)=\frac12\int_0^T\int_\Omega
       \bigl(|\partial_tu|^2-c^2|\nabla_xu|^2\bigr)\,\mathrm dx\,\mathrm dt.
\]
对在初、末时刻和空间边界均为零的扰动 \(h\)，时、空分别分部积分，得到
\[
 D\mathcal A(u)h=\int_0^T\int_\Omega
       (-\partial_t^2u+c^2\Delta_xu)h\,\mathrm dx\,\mathrm dt.
\]
所以驻点满足波动方程 \(\partial_t^2u-c^2\Delta u=0\)。这里固定初末值是推导变分方程时对扰动的限制；实际波动初值问题仍由初位移和初速度确定，不能把两种问题混为一谈。

\ProseParagraphBreak

作用量的二次型有正负两部分，一般并不在解处取得极小值。例如取乘积扰动 \(h(t,x)=a(t)b(x)\)，其中两因子均光滑紧支集，则
\[
 2\mathcal A(h)=\lVert a'\rVert_2^2\lVert b\rVert_2^2
      -c^2\lVert a\rVert_2^2\lVert\nabla b\rVert_2^2.
\]
固定非零 \(b\) 后，取 \(a_N(t)=\eta(t)\sin(Nt)\)，其中 \(0\ne\eta\in C_c^\infty(0,T)\)。展开正弦平方并用\cref{b2:thm:03-riemann-lebesgue}，有 \(\|a_N\|_2^2\to\|\eta\|_2^2/2\) 及 \(N^{-2}\|a_N'\|_2^2\to\|\eta\|_2^2/2\)，所以作用量为正。反之，固定非零 \(a\)，取 \(b_N(x)=\psi(x)\sin(Nx_1)\)，其中 \(0\ne\psi\in C_c^\infty(\Omega)\)。相同的平方展开给出 \(\|b_N\|_2^2\to\|\psi\|_2^2/2\) 及 \(N^{-2}\|\nabla b_N\|_2^2\to\|\psi\|_2^2/2\)，所以作用量为负。若 \(u\) 为驻点，精确二次展开给出 \(\mathcal A(u+th)-\mathcal A(u)=t^2\mathcal A(h)\)。故可把两类扰动都缩得任意小，在每个驻点附近仍有增加和降低作用量的方向。

\Needspace{6\baselineskip}
\begin{proposition}[波动能量守恒]
\label{b2:prop:14-wave-energy}
设整数 \(d\geq1\)，\(c,T>0\)，\(\Omega\subset\mathbb R^d\) 为非空有界 \(C^1\) 开集，实函数 \(u\) 在 \([0,T]\times\overline\Omega\) 的某开邻域内属于 \(C^2\)，满足 \(u_{tt}-c^2\Delta u=0\) 于 \((0,T)\times\Omega\)，且 \(u(t,y)=0\) 对 \(0\leq t\leq T\)、\(y\in\partial\Omega\) 成立。则
\[
 \mathcal E(t)=\frac12\int_\Omega
          \bigl(|u_t(t,x)|^2+c^2|\nabla u(t,x)|^2\bigr)\,\mathrm dx
\]
与时间无关。
\end{proposition}
\begin{proof}
对时间求导并对空间分部积分，得
\[
 \mathcal E'(t)=\int_\Omega u_t(u_{tt}-c^2\Delta u)\,\mathrm dx
           +c^2\int_{\partial\Omega}u_t\partial_\nu u\,\mathrm dS=0.
\]
内部项由方程为零，边界项由对恒定零边界条件求时间导数所得 \(u_t=0\) 而消失。
\end{proof}

守恒能量中空间项带正号，作用量中带负号。前者衡量一个时刻的状态，后者比较整段时空运动。对同样的区域、实值 \(C^2\) 解及恒定零边界，若方程改为 \(u_{tt}-c^2\Delta u+\gamma u_t=f\)，其中 \(\gamma\geq0\)、\(f\) 连续，则同一计算给出
\[
 \mathcal E'(t)=\int_\Omega f(t,x)u_t(t,x)\,\mathrm{d}x
                    -\gamma\|u_t(t)\|_2^2.
\]
这些等式将外界做功和耗散分开记录。\cref{b2:tab:14-action-energy}集中比较作用量与能量；表中的守恒结论始终包含零边界且无外力、无阻尼的条件。

\begin{table}[htbp]
\centering
\caption[波动作用量与守恒能量的区别]{波动作用量与守恒能量的区别。取整数 \(d\geq1\)、\(c,T>0\)、非空有界 \(C^1\) 开集 \(\Omega\subset\mathbb R^d\)，函数实值且光滑到足以完成正文的微分与分部积分，侧边界为恒定零值。作用量列的变分还固定初、末时刻的函数值；能量列讨论满足波动方程的实际演化。}
\label{b2:tab:14-action-energy}
\small
\begin{tabularx}{\textwidth}{@{}p{5.2em}XX@{}}
\toprule
比较项目 & 时空作用量 \(\mathcal A\) & 守恒能量 \(\mathcal E(t)\)\\
\midrule
积分范围 & 整段 \((0,T)\times\Omega\) & 固定时刻的 \(\Omega\)\\
空间梯度项 & \(-c^2|\nabla u|^2/2\) & \(+c^2|\nabla u|^2/2\)\\
数学作用 & 对允许扰动驻定，推出运动方程 & 控制速度与梯度，进而估计解的差\\
极值或演化 & 驻点一般为鞍点 & 无外力、无阻尼时随时间守恒\\
\bottomrule
\end{tabularx}
\end{table}

\subsection{PDE 与极值问题}
\label{b2:subsec:140303}

从微分方程寻找泛函时，最容易检验的是对称性。设 \(H\) 是实 Hilbert 空间，\(a:H\times H\to\mathbb R\) 是有界双线性型，\(\ell\in H^*\)。令
\[J(u)=\tfrac12a(u,u)-\ell(u).\]
直接展开可得
\[DJ(u)h=\tfrac12\bigl(a(u,h)+a(h,u)\bigr)-\ell(h).\]
因此仅当 \(a\) 对称时，驻点条件才是 \(a(u,h)=\ell(h)\)。反对称部分在 \(a(u,u)\) 中完全消失，却可能实质影响方程的解。

\ProseParagraphBreak

例如 \(H=\mathbb R^2\)，令 \(Au=(u_1-u_2,u_1+u_2)\)，\(a(u,h)=(Au)\cdot h\)。尽管 \(a(u,u)=|u|^2\) 严格为正，极小化 \(\tfrac12|u|^2-f\cdot u\) 得到的是 \(u=f\)，而原方程是 \(Au=f\)。两者一般不同。下一章的 Lax--Milgram 定理直接求解有界强制双线性型，无须先把它变成能量；这正是从变分极小化过渡到一般弱解方法时必须增加的能力。

\ProseParagraphBreak

另一个区别是驻点与动力学。若 \(E(u)=\tfrac12\int|\nabla u|^2\)，形式上的 \(L^2\) 梯度为 \(-\Delta u\)。沿最陡下降运动 \(u_t=-\operatorname{grad}E(u)\)，得到热方程 \(u_t=\Delta u\)，而平衡点满足 Laplace 方程。对光滑零边界解，分部积分给出
\[\frac{\mathrm d}{\mathrm dt}E(u(t))=-\int_\Omega|u_t|^2\,\mathrm dx\le0.\]
这项恒等式说明热演化趋向降低能量，却不在这里自动证明长时间收敛；还需要轨道紧性、极限点的识别等信息。变分结构可以提供估计和候选极限，但不能代替这些额外步骤。

% !TeX root = ../../Book2.tex
\section{直接方法}
\label{b2:sec:1404}

先推导 Euler--Lagrange 方程再求解，有时会把存在性问题变得更难。直接方法反过来考察能量本身：取能量趋近下确界的序列，控制它的大小，抽取极限，再证明极限没有丢失边界条件和能量下界。第一卷定理 9.5.9 已在自反 Banach 空间中证明这一抽象原则。本节不重复该证明，而逐项核验积分能量满足条件的具体原因。

\subsection{强制性}
\label{b2:subsec:140401}

在 Banach 空间 \(X\) 的约束集 \(K\) 上，泛函的强制性指
\[J(v)\longrightarrow+\infty\qquad(v\in K,\ \lVert v\rVert_X\to\infty).\]
它保证能量有界的候选者不能逃到范数无穷远处。在零边界空间 \(H_0^1(\Omega)\) 上，Poincaré 不等式和 Cauchy--Schwarz 不等式给出
\[
 \frac12\lVert\nabla v\rVert_2^2-\int_\Omega fv
 \ge\frac12\lVert\nabla v\rVert_2^2-C_P\lVert f\rVert_2\lVert\nabla v\rVert_2.
\]
右端是开口向上的二次式，故这项能量强制。若允许非零边界数据，写成 \(u=g+v\)，交叉项 \(\int\nabla g\cdot\nabla v\) 仍只按 \(\lVert\nabla v\rVert_2\) 线性增长，因此同样得到强制性。

\ProseParagraphBreak

边界限制不能从这个论证中删除。在 \(H^1(\Omega)\) 全空间中，常函数的梯度能量为零，故单靠梯度无法控制范数。若 \(\int_\Omega f\ne0\)，沿适当符号的常函数甚至能让总能量趋向 \(-\infty\)。Neumann 问题因此需要相容条件和一个固定常数的归一化；零均值空间上的 Poincaré 不等式才是相应替代物。

\subsection{极小化序列}
\label{b2:subsec:140402}

若 \(m=\inf_KJ\) 是有限实数，可以按下确界的定义选取 \(u_n\in K\)，满足 \(m\le J(u_n)<m+1/n\)。这样的序列称为极小化序列。选取过程既不要求驻点已存在，也不要求每个 \(u_n\) 近似满足 Euler--Lagrange 方程。

\ProseParagraphBreak

控制能量后，强制性使序列有界，但有界不等于强收敛。无限维 Hilbert 空间中的标准正交序列有界，却没有强收敛子列；它仅弱收敛到零。这就是直接方法引入弱拓扑的实质理由。另一方面，若约束集不是弱闭的，弱极限可能不再可行；若能量不是弱下半连续的，极限能量可能突然升高。序列有界、约束保留与能量比较分别承担不同任务，不能用一句“由紧性”全部略去。

\ProseParagraphBreak

二次能量还有一个可直接看见的特殊现象。若 \(u_n\) 和 \(u_k\) 满足同一边界数据，则
\[
 E\left(\frac{u_n+u_k}{2}\right)
 =\frac{E(u_n)+E(u_k)}2-\frac18\lVert\nabla(u_n-u_k)\rVert_2^2.
\]
左端不小于下确界 \(m\)。因此极小化序列事实上在梯度范数下为 Cauchy 序列，借完备性可以直接得到强极限。一般非二次能量不一定有这个精确恒等式，弱紧性的方法因而更具普遍性。

\subsection{弱紧性}
\label{b2:subsec:140403}

设 \(1<p<\infty\)。空间 \(W^{1,p}(\Omega)\) 可以等距嵌入 \((L^p(\Omega))^{d+1}\)，其像由弱导数关系确定，是闭子空间。有限个自反空间的乘积及其闭子空间仍自反；故第一卷定理 9.4.7 给出：每个有界 \(W^{1,p}\) 序列有弱收敛子列。具体而言，沿同一子列有
\[u_n\rightharpoonup u\ \text{于 }L^p,\qquad \partial_i u_n\rightharpoonup\partial_i u\ \text{于 }L^p.\]
弱导数关系可以直接验证：对测试函数 \(\varphi\)，在 \(\int u_n\partial_i\varphi=-\int(\partial_i u_n)\varphi\) 中取极限即可。

\ProseParagraphBreak

仿射边界类 \(g+W_0^{1,p}(\Omega)\) 是强闭凸集，因而弱闭。这里用的是第一卷定理 9.5.2 的闭凸集结论，完全不要求逐点读取边界值。若约束还包含 \(u\ge\psi\) 几乎处处，它也是强闭凸约束：强收敛子列可进一步逐点几乎处处收敛，从而保留不等式；闭凸性再给出弱闭性。

\ProseParagraphBreak

端点 \(p=1\) 与 \(p=\infty\) 不在上述自反性结论内。例如在 \((0,1)\) 上，\(u_n(x)=\min\{nx,1\}\) 的 \(W^{1,1}\) 范数有界，但导数集中到越来越短的边界层中，不能在 \(L^1\) 内取得弱收敛子列。其极限在内部是常数 \(1\)，左端数据却已丢失。处理线性增长能量时往往需要有界变差函数和测度导数；不能把 \(p>1\) 的论证直接搬到这个端点。

\subsection{弱下半连续性}
\label{b2:subsec:140404}

第一卷定理 9.5.8 说明：凸泛函的范数下半连续性等价于弱下半连续性。对积分泛函，一个常用核验顺序是先证强收敛下的下半连续，再证凸性。设 \(F(x,z,p)\ge0\) 对 \((z,p)\) 连续且凸，对 \(x\) 可测，则
\[I(u)=\int_\Omega F(x,u,\nabla u)\,\mathrm dx\]
是允许取 \(+\infty\) 的凸泛函。若 \(u_n\to u\) 于 \(W^{1,p}\)，先抽取实现能量下极限的子列，再从中抽取函数与全部一阶导数几乎处处收敛的子列。连续性与 Fatou 引理给出 \(I(u)\le\liminf I(u_n)\)。因此 \(I\) 强下半连续，继而弱下半连续。

\ProseParagraphBreak

非负性可放宽为统一的可积下界，因为减去该下界不改变论证。又若 \(\ell\in(W^{1,p})^*\)，线性项在弱收敛下连续，所以 \(I-\ell\) 仍弱下半连续。这个证明覆盖了 \(\tfrac1p\int|\nabla u|^p\) 及许多非负凸势能，却没有覆盖任意非凸密度。特别是向量值问题中，关于梯度的普通凸性虽为充分条件，却远非最一般的下半连续性条件；本章不把更深的条件隐含在这里。

\ProseParagraphBreak

以 \(\ell^2\) 上的 \(J(x)=-\lVert x\rVert^2\) 为例，\(J\) 范数连续，但对标准正交序列 \(e_n\rightharpoonup0\) 有 \(J(0)=0>-1=\lim J(e_n)\)。所以范数连续本身不能代替弱下半连续；前面的凸性在改变拓扑时起了决定作用。

\subsection{极小元的存在性}
\label{b2:subsec:140405}

下面把这些核验汇集到一个非线性模型。它既包含 Dirichlet 能量，也说明直接方法并不依赖显式解或线性叠加。

\begin{theorem}[梯度幂能量的极小元]
\label{b2:thm:14-p-energy-minimum}
设 \(\Omega\subset\mathbb R^d\) 是有界开集，\(1<p<\infty\)，\(g\in W^{1,p}(\Omega)\)，\(f\in L^{p'}(\Omega)\)，其中 \(p'=p/(p-1)\)。则
\[
 J(u)=\frac1p\int_\Omega|\nabla u|^p\,\mathrm dx-\int_\Omega fu\,\mathrm dx
\]
在 \(g+W_0^{1,p}(\Omega)\) 上有唯一极小元 \(u\)。它满足
\[
 \int_\Omega|\nabla u|^{p-2}\nabla u\cdot\nabla h\,\mathrm dx
     =\int_\Omega fh\,\mathrm dx\qquad(h\in W_0^{1,p}(\Omega)),
\]
其中在 \(\nabla u=0\) 处约定 \(|\nabla u|^{p-2}\nabla u=0\)。
\end{theorem}
\begin{proof}
令 \(u=g+v\)。由三角不等式和 Poincaré 不等式，\(\lVert\nabla v\rVert_p\le\lVert\nabla u\rVert_p+\lVert\nabla g\rVert_p\)，且
\[\left|\int_\Omega fu\right|\le\lVert f\rVert_{p'}\bigl(\lVert g\rVert_p+C_P\lVert\nabla v\rVert_p\bigr).\]
因此 \(J(g+v)\) 的下界由一个正的 \(\lVert\nabla u\rVert_p^p\) 项减去线性项构成。当 \(\lVert v\rVert_{W^{1,p}}\to\infty\) 时，\(\lVert\nabla u\rVert_p\to\infty\)，而 \(p>1\)，故泛函强制且有有限下界。可行集非空、弱闭；上一小节已核验梯度幂积分弱下半连续，而 \(f\) 的配对弱连续。于是自反性及第一卷定理 9.5.9 给出极小元。

若 \(u_1,u_2\) 都极小，由 \(\xi\mapsto|\xi|^p\) 的严格凸性，除非 \(\nabla u_1=\nabla u_2\) 几乎处处，中点能量就严格更小。故梯度相等；差属于 \(W_0^{1,p}\)，Poincaré 不等式推出差为零。

最后固定 \(h\in W_0^{1,p}\)。沿线段求导时，\(|t|\le1\) 的导数被
\[C_p\bigl(|\nabla u|^{p-1}+|\nabla h|^{p-1}\bigr)|\nabla h|\]
支配；该函数由 Hölder 不等式可积。支配收敛允许交换求导与积分，得第一变分公式。极小点的正负方向都可行，故第一变分为零，即所列等式。
\end{proof}

当 \(p=2\) 时，这就是 Poisson 方程的弱形式。一般 \(p\) 则对应 \(-\operatorname{div}(|\nabla u|^{p-2}\nabla u)=f\)，通常称为 \(p\)-Laplace 方程。上述存在性并没有证明极小元是 \(C^2\) 函数；能写出弱方程与能把它逐点展开，是两个不同层次的结论。

\ProseParagraphBreak

如果再加上非负连续凸势能 \(\int_\Omega\Psi(u)\)，只要某个可行函数具有有限能量，同样的直接方法仍给出极小元；严格凸的梯度项保留唯一性。但要写出 \(\Psi'(u)\) 的弱方程，还须补充可微性和增长条件，确保扰动可积并允许积分下求导。存在性所需的下半连续性通常比推导经典微分方程所需的光滑性弱得多，这也是直接方法的重要优势。

% 本章小结（不计入编号）
\begin{chaptersummary}
第一变分把极小问题转成对全部允许方向成立的积分等式；分部积分把该等式拆成内部 Euler--Lagrange 方程和自然边界条件。允许的方向由约束决定，故固定边界、自由边界及单侧约束产生不同形式。第二变分非负只是必要条件，一致正下界才给出局部极小性的稳定判据。对凸二次能量，精确能量差进一步把驻点与全局唯一极小点联系起来。

\ProseParagraphBreak

直接方法不先求解微分方程，而从极小化序列出发。强制性控制序列，Sobolev 空间的自反性提供弱极限，闭凸约束保留边界数据，弱下半连续性完成能量比较。所获极小元首先属于弱函数空间；它是否更光滑，需要另外的方程估计。下一章将以弱形式直接求解一般双线性问题，尤其处理不能来自对称能量的算子。
\end{chaptersummary}


\BookExercises

下面各题均为本书自拟。题目通过改变空间、边界或能量符号，检验正文结论的适用范围。

% 自拟：无限维正定第二变分不足以保证极小。
\begin{exercise}\label{b2:ex:14-positive-hessian}
在实空间 \(\ell^2\) 上令 \(J(x)=\sum_{n\ge1}x_n^2/n-\lVert x\rVert^4\)。证明 \(J\) 二次连续 Fréchet 可微，零点为驻点且 \(D^2J(0)[h,h]>0\) 对每个 \(h\ne0\) 成立，但零点不是局部极小点。指出第二变分判别的哪项条件不成立。
\end{exercise}
\begin{solution}
有界自伴算子 \(Ax=(x_n/n)\) 给出 \(J(x)=\langle Ax,x\rangle-\langle x,x\rangle^2\)，所以它是连续多项式，特别是二次连续可微。\(DJ(0)=0\)，且 \(D^2J(0)[h,h]=2\sum h_n^2/n>0\)。令 \(x^{(n)}=\sqrt{2/n}\,e_n\)，则 \(\lVert x^{(n)}\rVert\to0\)，而 \(J(x^{(n)})=2/n^2-4/n^2=-2/n^2<0=J(0)\)。在单位向量 \(e_n\) 上第二变分等于 \(2/n\to0\)，故不存在统一正下界 \(c\lVert h\rVert^2\)。
\end{solution}

% 自拟：带权端点问题，应用驻点和严格凸性。
\begin{exercise}\label{b2:ex:14-weighted-string}
设 \(a\in C^1([0,1])\)、\(a\ge a_0>0\)。求 \(\tfrac12\int_0^1a(x)|u'(x)|^2\,\mathrm dx\) 在 \(u(0)=0,u(1)=1\) 的 \(H^1\) 函数中的极小元和极小值，并用能量差证明唯一性。
\end{exercise}
\begin{solution}
记 \(L=\int_0^1a(s)^{-1}\,\mathrm ds\)，令 \(u(x)=L^{-1}\int_0^x a(s)^{-1}\,\mathrm ds\)，则 \(au'=L^{-1}\)。任一零端点扰动 \(h\in H_0^1(0,1)\) 满足 \(\int au'h'=L^{-1}(h(1)-h(0))=0\)，此等式也可由光滑逼近得到。因此能量差为 \(\tfrac12\int a|h'|^2\ge0\)，且等号迫使 \(h'=0\)、\(h=0\)。极小值是 \((2L^2)^{-1}\int a^{-1}=1/(2L)\)。
\end{solution}

% 自拟：边界弹簧与外法向符号。
\begin{exercise}\label{b2:ex:14-boundary-spring}
给定 \(\beta>0\)、\(q\in\mathbb R\)，在 \(u\in H^1(0,1)\)、\(u(0)=0\) 中极小化
\[J(u)=\tfrac12\int_0^1|u'|^2\,\mathrm dx+\tfrac\beta2u(1)^2-qu(1).\]
求极小元，并核验自然边界条件与全局极小性。
\end{exercise}
\begin{solution}
第一变分为 \(\int u'h'+(\beta u(1)-q)h(1)\)，其中 \(h(0)=0\)。内部方程是 \(u''=0\)，右端条件是 \(u'(1)+\beta u(1)=q\)，所以 \(u(x)=qx/(1+\beta)\)。对任意允许扰动，交叉项消失，能量差为 \(\tfrac12\int|h'|^2+\tfrac\beta2h(1)^2\)。它非负，等号时 \(h=0\)，故所得解是唯一全局极小元。
\end{solution}

% 自拟：积分约束产生常数乘子，含显式解。
\begin{exercise}\label{b2:ex:14-mass-constraint}
给定 \(M\in\mathbb R\)，在 \(H_0^1(0,1)\) 中，附加 \(\int_0^1u=M\)，极小化 \(\tfrac12\int_0^1|u'|^2\)。不用抽象乘子定理，求出唯一极小元和极小值。
\end{exercise}
\begin{solution}
取 \(u(x)=6Mx(1-x)\)，则端点为零且积分为 \(M\)。允许扰动 \(h\) 满足零端点和零积分；分部积分给出 \(\int u'h'=-\int u''h=12M\int h=0\)。能量差为 \(\tfrac12\lVert h'\rVert_2^2\)，故它是唯一极小元。又 \(\int_0^1(1-2x)^2\,\mathrm dx=1/3\)，极小值为 \(6M^2\)。常数 \(-u''=12M\) 正是积分约束所允许留下的乘子；若错误地令所有紧支集方向都可行，就会误得 \(u''=0\)。
\end{solution}

% 自拟：临界负势能与最优Poincare常数。
\begin{exercise}\label{b2:ex:14-negative-potential}
在 \(H_0^1(0,\pi)\) 上考察 \(J_\lambda(u)=\tfrac12\int_0^\pi(|u'|^2-\lambda|u|^2)\)。使用区间正弦展开所得 \(\lVert u\rVert_2\le\lVert u'\rVert_2\)，分别讨论 \(\lambda<1\)、\(\lambda=1\)、\(\lambda>1\) 的极小元与强制性。
\end{exercise}
\begin{solution}
若 \(\lambda<1\)，当 \(0\le\lambda<1\) 时 \(J_\lambda\ge(1-\lambda)\lVert u'\rVert_2^2/2\)，当 \(\lambda<0\) 时更有 \(J_\lambda\ge\lVert u'\rVert_2^2/2\)，故强制且唯一极小元为零。若 \(\lambda=1\)，正弦展开给出 \(2J_1=\frac\pi2\sum_{n\ge1}(n^2-1)|a_n|^2\)，所以极小元恰为 \(c\sin x\)，泛函不强制。若 \(\lambda>1\)，沿 \(u=c\sin x\) 有 \(J_\lambda=\pi(1-\lambda)c^2/4\to-\infty\)，故没有有限极小值。展开公式先对有限正弦和成立，再由 \(H_0^1\) 逼近推广。
\end{solution}

% 自拟：非对称双线性型的有限维反例。
\begin{exercise}\label{b2:ex:14-skew-part}
令 \(A=\left(\begin{smallmatrix}2&-3\\3&2\end{smallmatrix}\right)\)、\(f=(1,0)\)。分别求 \(Au=f\) 的解及 \(\tfrac12u^{\mathsf T}Au-f\cdot u\) 的极小元。证明 \((Au)\cdot u\ge2|u|^2\)，并解释这一强制性为何不足以使两个问题相同。
\end{exercise}
\begin{solution}
\(A^{-1}=13^{-1}\left(\begin{smallmatrix}2&3\\-3&2\end{smallmatrix}\right)\)，故方程解为 \((2/13,-3/13)\)。由于 \(u^{\mathsf T}Au=2|u|^2\)，泛函是 \(|u|^2-u_1\)，唯一极小元为 \((1/2,0)\)。强制性等式确实成立，但二次型只保留 \((A+A^{\mathsf T})/2=2I\)；驻点方程为 \(2u=f\)，反对称部分没有进入能量。
\end{solution}

% 自拟：障碍约束产生变分不等式。
\begin{exercise}\label{b2:ex:14-obstacle}
设 \(K\) 是 \(H_0^1(\Omega)\) 的非空闭凸集，\(f\in L^2(\Omega)\)，\(\Omega\) 有界。证明 Dirichlet 外力能量在 \(K\) 上有唯一极小元 \(u\)，并证明它等价于
\[\int_\Omega\nabla u\cdot\nabla(v-u)\ge\int_\Omega f(v-u)\qquad(v\in K).\]
\end{exercise}
\begin{solution}
Poincaré 不等式给出强制性；闭凸集弱闭，能量凸且范数连续，故弱下半连续。第一卷定理 9.5.9 给出存在性，梯度平方的严格凸性给出唯一性。若 \(u\) 极小，对 \(u+t(v-u)\in K\)、\(0<t\le1\) 展开能量，除以 \(t\) 令 \(t\downarrow0\)，即得不等式。反之，取 \(h=v-u\) 于精确二次展开，线性项由不等式非负，平方项也非负，故 \(E(v)\ge E(u)\)。凸约束只提供单侧可行线段，因而得到不等式。
\end{solution}

% 自拟：受迫阻尼波的能量估计。
\begin{exercise}\label{b2:ex:14-damped-wave}
设光滑零边界函数满足 \(u_{tt}-c^2\Delta u+\gamma u_t=f\)，其中 \(\gamma>0\)。证明对 \(0\le t\le T\)，
\[\mathcal E(t)+\frac\gamma2\int_0^t\lVert u_s\rVert_2^2\,\mathrm ds
 \le\mathcal E(0)+\frac1{2\gamma}\int_0^t\lVert f(s)\rVert_2^2\,\mathrm ds.\]
说明该估计怎样给出相同初边值数据的唯一性。
\end{exercise}
\begin{solution}
乘方程以 \(u_t\) 并分部积分，得 \(\mathcal E'+\gamma\lVert u_t\rVert_2^2=\int fu_t\)。由 Young 不等式，右端不超过 \(\lVert f\rVert_2^2/(2\gamma)+\gamma\lVert u_t\rVert_2^2/2\)。移项并对时间积分，得到所需估计。对两个解之差应用，右端外力与初始能量都为零，故差的时间导数和梯度均为零；零边界或零初位移固定剩余常数，差恒为零。
\end{solution}

% !TeX root = ../../Book2.tex
\chapter{弱解与 Lax–Milgram 方法}
\label{b2:ch:15}
\BookChapterNumber{b2:ch:15}
% 本章摘要（不计入编号）
\begin{chapterabstract}
本章把散度型方程写成 Hilbert 空间中的测试等式。先明确 Sobolev 边界条件与负阶源项，再完整证明 Lax--Milgram 定理，并以 Poisson 方程说明存在性、唯一性、稳定性和有限维近似。一般有界可测系数无需具有经典导数；一致椭圆性控制主部，Gårding 不等式处理低阶项。最后借紧嵌入与第一卷的 Fredholm 理论，刻画原方程可能出现的有限维核及伴随相容条件。
\end{chapterabstract}

% 本章目标（不计入编号）
\begin{learninggoals}
\begin{itemize}
\item\label{b2:item:15-goal-weak}辨认测试空间所编码的内部方程与边界条件，正确使用 \(H^{-1}\) 数据。
\item\label{b2:item:15-goal-lax}分别核验型的有界性、强制性与空间完备性，掌握 Lax--Milgram 的闭值域证明。
\item\label{b2:item:15-goal-stability}用能量等式证明唯一性、比较和数据稳定性，并处理非齐次边界延拓。
\item\label{b2:item:15-goal-general}说明有界可测系数已足以给出弱存在性，区别原问题与移位问题。
\item\label{b2:item:15-goal-fredholm}在选读范围内，用伴随核写出 Fredholm 相容条件。
\end{itemize}
\end{learninggoals}

设一根杆的导热系数在接头处突然变化。平衡温度不应跳跃，穿过接头的热通量也应保持一致，但温度斜率可以不同。若只寻找处处二次可微的解，便会把这样的基本模型排除在外。弱形式保留可积梯度和通量平衡，把经典导数无法直接表达的问题放进一个完备的函数空间。

\ProseParagraphBreak

上一章用极小化构造对称能量的解；本章直接求解双线性等式，从而允许非对称系数。Hilbert 空间表示与紧算子的抽象结论采用第一卷第八章和附录 C 的证明，Sobolev 边界与紧嵌入采用第一卷第二十二至二十五章。相关理论的系统参考包括 Haim Brezis 的《Functional Analysis, Sobolev Spaces and Partial Differential Equations》\cite{Brezis2011Functional}与 Lawrence C. Evans 的《Partial Differential Equations》\cite{Evans2010Partial}。

% !TeX root = ../../Book2.tex
\section{弱形式}
\label{b2:sec:1501}

分段材料中的温度可以连续而斜率发生跳跃；此时逐点写二阶导数会掩盖界面上的通量条件。弱形式把微分移到测试函数上，使未知函数只需具有有限能量，同时仍保留方程及界面平衡。本节将第一章的局部弱解概念与上一章的 Sobolev 边界类结合起来，并为比 \(L^2\) 更粗糙的源项确定合适空间。

\subsection{测试函数}
\label{b2:subsec:150101}

在开集 \(\Omega\subset\mathbb R^d\) 中，测试函数取自 \(C_c^\infty(\Omega)\)。它们的支集离边界有正距离，因此分部积分不会引入尚未定义的边界值。若 \(u\in H_{\mathrm{loc}}^1(\Omega)\)、\(A\in L_{\mathrm{loc}}^\infty(\Omega;\mathbb R^{d\times d})\)，则
\[\langle-\operatorname{div}(A\nabla u),\varphi\rangle
 =\int_\Omega A\nabla u\cdot\nabla\varphi\,\mathrm dx\]
定义了一个分布。矩阵约定为 \((A\nabla u)_i=\sum_j a_{ij}\partial_j u\)，因而散度作用于通量的第一个指标。

\ProseParagraphBreak

局部测试函数足以辨认内部方程，但能量估计常须取 \(\varphi=u\)、截断后的 \(u\) 或差商。它们通常不光滑，也不紧支集。合法性依靠的是密度与连续延拓：若左右两边对 \(H_0^1\) 范数连续，而 \(C_c^\infty\) 在该空间稠密，则测试等式自动推广到整个 \(H_0^1\)。若只知道局部 \(H^1\)，便先乘紧支集截断函数，在较小区域完成同样步骤。不能在未建立这些连续性之前直接“取测试函数为解”。

\subsection{弱解}
\label{b2:subsec:150102}

以下主要讨论实函数，复函数可分别作用于实部与虚部。给定有界矩阵场 \(A\) 和源项 \(f\)，称 \(u\in H^1(\Omega)\) 满足散度方程的弱形式，是指
\[\int_\Omega A\nabla u\cdot\nabla\varphi=\langle f,\varphi\rangle
 \qquad(\varphi\in C_c^\infty(\Omega)).\]
这与第一章的定义一致；边界条件尚未包含在此等式中。如果 \(A,u\) 光滑到足以分部积分，它与经典方程等价；粗糙情形则以这个等式本身为准。

\begin{proposition}[一维界面的通量条件]
\label{b2:prop:15-interface-flux}
设 \(-1<s<1\)，\(a\) 在 \((-1,s)\)、\((s,1)\) 上分别为正常数，\(u\) 在两闭子区间上分别为 \(C^2\)，且在 \(s\) 处连续。若 \(-(au')'=f\) 在两侧经典成立，其中 \(f\in L^2(-1,1)\)，则 \(u\) 在整个区间为弱解，当且仅当
\[a(s-)u'(s-)=a(s+)u'(s+).\]
\end{proposition}
\begin{proof}
分段光滑且连续的 \(u\) 属于 \(H^1\)，其弱导数就是分段的一阶导数。在两段分别积分，外端因测试函数紧支集而消失，内端得到
\[\int_{-1}^1 au'\varphi'=\int_{-1}^1 f\varphi
  +\bigl(a(s-)u'(s-)-a(s+)u'(s+)\bigr)\varphi(s).\]
测试函数在 \(s\) 处可取任意值，因此等式对全部测试函数成立恰当且仅当括号为零。
\end{proof}

例如 \(a=1\) 于左半区间、\(a=2\) 于右半区间，令 \(u(x)=x\)（\(x\le0\)）、\(u(x)=x/2\)（\(x\ge0\)）。其导数跳跃，但通量恒为 \(1\)，所以 \(-(au')'=0\) 弱成立。若误要求 \(u'\) 连续，就会排除这个正确的平衡解；若允许 \(u\) 本身跳跃，又会产生分布导数中的点质量，通常已不属于所选能量空间。

\subsection{Sobolev 边界条件}
\label{b2:subsec:150103}

对任意有界开集，齐次 Dirichlet 条件定义为 \(u\in H_0^1(\Omega)\)；非齐次数据由一个延拓 \(g\in H^1(\Omega)\) 表示，条件为 \(u-g\in H_0^1(\Omega)\)。第一卷定理 23.1.6 说明，在有界 Lipschitz 区域上，这恰好等价于 Sobolev 迹满足 \(\operatorname{Tr}u=\operatorname{Tr}g\)。粗糙区域上仍可使用闭包定义，但不能未经证明把它等同于某个朴素逐点边界条件。

\ProseParagraphBreak

自然边界条件的写法不同。以 Poisson 方程为例，在 \(H^1(\Omega)\) 的全部方向测试时，光滑情形的 Green 公式是
\[\int_\Omega\nabla u\cdot\nabla v
  =\int_\Omega fv+\int_{\partial\Omega}g_N\operatorname{Tr}v\,\mathrm dS.\]
它编码 \(\partial_\nu u=g_N\)，而不是要求 \(u\) 在边界为零。取 \(v=1\) 立刻得到相容条件 \(\int_\Omega f+\int_{\partial\Omega}g_N=0\)。测试空间不同，编码的边界问题便不同；仅看内部方程无法辨认二者。

\ProseParagraphBreak

即使区域光滑，也不宜把任意边界函数都当作合法 \(H^1\) Dirichlet 数据。数据应当处于迹空间，并存在 \(H^1\) 延拓；第一卷第二十三章在 Lipschitz 图中用分数阶范数刻画这种要求。本章直接给定延拓 \(g\)，从而把数据的合法性写进假设，不需要在每个存在性证明里重新构造延拓算子。

\subsection{负阶 Sobolev 空间初步}
\label{b2:subsec:150104}

设 \(\Omega\) 有界，以梯度范数赋予 \(H_0^1(\Omega)\) Hilbert 结构。定义
\[H^{-1}(\Omega)=(H_0^1(\Omega))^*,\qquad
 \lVert f\rVert_{H^{-1}}=\sup_{\lVert\nabla v\rVert_2\le1}|\langle f,v\rangle|.
\]
这是本章使用的 \term[fujie-sobolev-kongjian]{负阶 Sobolev 空间}[negative-order Sobolev space]的第一种情形。换成全 \(H^1\) 范数会得到等价范数；双空间本身不变。由 Poincaré 不等式，\(L^2\) 连续嵌入 \(H^{-1}\)。更一般地，若 \(F\in L^2(\Omega;\mathbb R^d)\)，则 \(-\operatorname{div}F\in H^{-1}\)，且其范数不超过 \(\lVert F\rVert_2\)。

\begin{proposition}[负阶数据的通量表示]
\label{b2:prop:15-negative-representation}
对每个 \(f\in H^{-1}(\Omega)\)，存在 \(F\in L^2(\Omega;\mathbb R^d)\)，使
\[\langle f,v\rangle=\int_\Omega F\cdot\nabla v\,\mathrm dx
 \quad(v\in H_0^1(\Omega)),\qquad \lVert F\rVert_2=\lVert f\rVert_{H^{-1}}.\]
\end{proposition}
\begin{proof}
映射 \(v\mapsto\nabla v\) 把 \(H_0^1\) 等距嵌入 \(L^2(\Omega;\mathbb R^d)\)，其像因完备性而闭。在线性子空间上定义 \(\Lambda(\nabla v)=\langle f,v\rangle\)；Poincaré 不等式保证定义无歧义。Hahn--Banach 定理将其等范数延拓到整个向量值 \(L^2\)。第一卷定理 8.4.7 的 Hilbert 空间 Riesz 表示给出向量 \(F\)，其范数等于泛函范数，恰得所列公式。
\end{proof}

这种表示未必唯一，因为与所有梯度正交的通量可以自由加到 \(F\) 上。负阶空间描述的是通量的散度，而不是某个唯一的通量。它也确实容纳比函数更粗糙的数据：在一维有界区间内，点值由 \(H_0^1\) 范数控制，故点质量属于 \(H^{-1}\)；在二维及更高维，一般不再如此。后面的习题用缩放说明这一维数差异，避免把一维直觉直接移到高维。

% !TeX root = ../../Book2.tex
\section{Hilbert 空间方法}
\label{b2:sec:1502}

上一章的极小化要求能量保留方程的全部信息。对非对称系数，这个要求可能失败；双线性型方法则直接研究等式 \(a(u,v)=\ell(v)\)。其中未知量占第一个位置，测试向量占第二个位置。本节把有界性和强制性分别转成算子的上界和下界，再用正交性证明可解性。

\subsection{双线性型}
\label{b2:subsec:150201}

设 \(H\) 是实 Hilbert 空间。\term[shuangxianxing-xing]{双线性型}[bilinear form]是对两个变量分别线性的映射 \(a:H\times H\to\mathbb R\)。对称性 \(a(u,v)=a(v,u)\) 是额外要求，本章的主要存在定理不作这一假设。例如
\[a(u,v)=\int_\Omega A\nabla u\cdot\nabla v\,\mathrm dx\]
在 \(H_0^1(\Omega)\) 上是双线性型；若矩阵 \(A(x)\) 对称，则型对称。一个纯反对称矩阵场对 \(a(u,u)\) 没有贡献，却可能改变 \(a(u,v)\)。

\ProseParagraphBreak

复 Hilbert 空间中通常改用对第一个变量线性、对第二个变量共轭线性的半双线性型。相应右端也对测试向量共轭线性，强制性写为 \(\operatorname{Re}a(u,u)\ge\alpha\lVert u\rVert^2\)。这与本书内积对第一变量线性的约定一致。为使证明中的对偶方向透明，下面先给出实空间版本，随后说明复版本的调整；不要把这一约定与分布配对的双线性约定混用。

\subsection{有界性}
\label{b2:subsec:150202}

双线性型的有界性是存在 \(M<\infty\)，使
\[|a(u,v)|\le M\lVert u\rVert_H\lVert v\rVert_H\qquad(u,v\in H).\]
它保证固定任一变量后所得泛函连续，也保证从稠密的测试子空间推广等式时极限合法。若 \(A\) 的矩阵算子范数几乎处处不超过 \(\Lambda\)，则
\[|a(u,v)|\le\Lambda\lVert\nabla u\rVert_2\lVert\nabla v\rVert_2.\]
这里有界性只使用系数的本性上界，不涉及其导数，也不要求点点连续。

\begin{proposition}[有界型的算子表示]
\label{b2:prop:15-form-operator}
设 \(a\) 是实 Hilbert 空间 \(H\) 上的有界双线性型。则存在唯一有界线性算子 \(A:H\to H\)，使
\[a(u,v)=\langle Au,v\rangle_H\qquad(u,v\in H),\]
且 \(\lVert A\rVert=\sup_{\lVert u\rVert,\lVert v\rVert\le1}|a(u,v)|\)。型对称当且仅当算子自伴。
\end{proposition}
\begin{proof}
固定 \(u\)，第一卷定理 8.4.7 把连续线性泛函 \(v\mapsto a(u,v)\) 表示为与唯一向量 \(Au\) 的内积。表示的唯一性和第一变量线性性给出 \(A\) 的线性性，且 \(\lVert Au\rVert\le M\lVert u\rVert\)。对单位向量分别取上确界给出范数等式。最后，\(a(u,v)=a(v,u)\) 等价于 \(\langle Au,v\rangle=\langle u,Av\rangle\)，正是自伴条件。
\end{proof}

这个 \(A\) 是依赖所选 Hilbert 内积的表示算子，不必等于微分表达式本身。例如在梯度内积的 \(H_0^1\) 中，Poisson 型 \(\int\nabla u\cdot\nabla v\) 的表示算子就是恒等算子；作为分布的 \(-\Delta\) 则把 \(H_0^1\) 映到 \(H^{-1}\)。分清值域有助于避免把有界算子与无界微分算子的谱性质混在一起。

\subsection{强制性}
\label{b2:subsec:150203}

称型 \(a\) 强制，是指存在 \(\alpha>0\)，使
\[a(u,u)\ge\alpha\lVert u\rVert_H^2\qquad(u\in H).\]
这里要求统一的 \(\alpha\)，而不只是每个非零向量上取正值。结合算子表示与 Cauchy--Schwarz 不等式，立刻得到
\[\alpha\lVert u\rVert_H\le\lVert Au\rVert_H.\]
它同时给出单射性和逆映射在值域上的有界性。不过，单有这种下界并不保证一般算子满射；例如 \(\ell^2\) 上的右移是等距单射却漏掉第一个坐标方向。Lax--Milgram 证明还要利用同一个型在测试向量处的正性，排除值域的正交补。

\ProseParagraphBreak

对椭圆型矩阵，条件 \(\xi^{\mathsf T}A(x)\xi\ge\lambda|\xi|^2\) 几乎处处成立时，梯度型在 \(H_0^1\) 的梯度范数下以 \(\alpha=\lambda\) 强制。在整个 \(H^1\) 上却并非如此，因为常函数在型中为零。选空间与选范数是定理假设的一部分，不能在写下“椭圆”之后省略核验。

\subsection{Lax–Milgram 定理}
\label{b2:subsec:150204}

\begin{theorem}[Lax--Milgram]
\label{b2:thm:15-lax-milgram}
设 \(H\) 是实 Hilbert 空间，\(a:H\times H\to\mathbb R\) 是双线性型，且对某些 \(M<\infty\)、\(\alpha>0\) 有
\[|a(u,v)|\le M\lVert u\rVert_H\lVert v\rVert_H,
 \qquad a(u,u)\ge\alpha\lVert u\rVert_H^2.\]
则对每个 \(\ell\in H^*\)，存在唯一 \(u\in H\) 满足
\[a(u,v)=\ell(v)\qquad(v\in H),\]
且 \(\lVert u\rVert_H\le\alpha^{-1}\lVert\ell\rVert_{H^*}\)。
\end{theorem}
\begin{proof}
用 \cref{b2:prop:15-form-operator} 表示 \(a(u,v)=\langle Au,v\rangle\)，并用 Riesz 表示 \(\ell(v)=\langle f,v\rangle\)。所求等式等价于 \(Au=f\)。强制性给出 \(\lVert Au\rVert\ge\alpha\lVert u\rVert\)，所以 \(A\) 单射。

若 \(Au_n\to y\)，下界说明 \((u_n)\) 是 Cauchy 序列，因而收敛到某个 \(u\in H\)。由 \(A\) 连续，\(Au=y\)，故值域闭。若 \(w\) 与值域正交，则 \(a(v,w)=\langle Av,w\rangle=0\) 对所有 \(v\in H\) 成立。取 \(v=w\)，强制性推出 \(w=0\)。因此值域的正交补为零；闭子空间的正交分解给出值域等于 \(H\)。于是 \(Au=f\) 唯一可解，且算子下界给出 \(\lVert u\rVert\le\alpha^{-1}\lVert f\rVert=\alpha^{-1}\lVert\ell\rVert\)。
\end{proof}

定理也称 \term[lax-milgram-dingli]{Lax--Milgram 定理}[Lax--Milgram theorem]。复空间版本把强制性改成实部下界，并把 \(\ell\) 视为连续共轭线性泛函；算子表示、闭值域和正交补的三步证明保持不变，在取 \(v=w\) 时使用 \(\operatorname{Re}a(w,w)\) 即可。

\ProseParagraphBreak

还可以把存在性看成一次收缩迭代。若 \(A\) 满足上述上下界，取 \(0<t<2\alpha/M^2\)，则
\[
 \lVert(I-tA)w\rVert^2
 =\lVert w\rVert^2-2t\langle Aw,w\rangle+t^2\lVert Aw\rVert^2
 \le(1-2t\alpha+t^2M^2)\lVert w\rVert^2.
\]
因此 \(u_{n+1}=u_n+t(f-Au_n)\) 收敛到唯一解。这个计算不仅提供另一种构造，也表明当 \(\alpha\) 很小而 \(M\) 很大时，允许的步长和收敛速度都会变差。定理的稳定性常数真实反映了问题接近失去强制性的程度。

% !TeX root = ../../Book2.tex
\section{Poisson 方程}
\label{b2:sec:1503}

现在把抽象定理用于最基本的边值问题。区域只需有界，源项只需属于 \(H^{-1}\)；无需显式 Green 函数，也无需预先知道解的二阶导数。存在性、稳定性与边界数据的处理都在同一个能量空间中完成。

\subsection{弱解存在性}
\label{b2:subsec:150301}

\begin{theorem}[Dirichlet 问题的弱可解性]
\label{b2:thm:15-poisson-existence}
设 \(\Omega\subset\mathbb R^d\) 为有界开集，\(f\in H^{-1}(\Omega)\)。则存在唯一 \(u\in H_0^1(\Omega)\)，使
\[\int_\Omega\nabla u\cdot\nabla v\,\mathrm dx=\langle f,v\rangle
 \qquad(v\in H_0^1(\Omega)).\]
并且 \(\lVert\nabla u\rVert_2\le\lVert f\rVert_{H^{-1}}\)。
\end{theorem}
\begin{proof}
取 \(H=H_0^1(\Omega)\)，内积为 \(\langle u,v\rangle_H=\int\nabla u\cdot\nabla v\)。由 Poincaré 不等式，它与完整 Sobolev 范数等价，故 \(H\) 完备。型 \(a(u,v)=\langle u,v\rangle_H\) 的有界常数与强制常数均为 \(1\)，右端按 \(H^{-1}\) 定义连续。\cref{b2:thm:15-lax-milgram} 因而给出存在、唯一性及估计。
\end{proof}

本例中 Lax--Milgram 退化为 Riesz 表示；一般系数时，同一个定理才显出更大的作用。该弱解也恰为 \(\tfrac12\lVert\nabla v\rVert_2^2-\langle f,v\rangle\) 的唯一极小元：只要将上一章的二次展开中的 \(L^2\) 配对换成对偶配对，证明逐字成立。两种构造在对称情形相遇，但弱型方法还可以继续处理非对称问题。

\ProseParagraphBreak

对 \(\Omega=(0,1)\)、\(f=\delta_s\)、\(0<s<1\)，解为
\[
 u(x)=\begin{cases}(1-s)x,&0\le x\le s,\\s(1-x),&s\le x\le1.\end{cases}
\]
它连续、分段仿射，导数在 \(s\) 处向下跳 \(1\)。两段分部积分给出 \(\int_0^1u'v'=v(s)\)，先对光滑 \(v\) 成立，再由点值的连续性推广。因此点源产生一个合法弱解，却不产生经典 \(C^2\) 解；这不是存在性定理的缺陷，而是数据本身的粗糙性。

\subsection{唯一性}
\label{b2:subsec:150302}

唯一性的能量论证值得单独保留，因为它不依赖解是如何构造的。若 \(u_1,u_2\) 具有同一数据，差 \(w\in H_0^1\) 满足 \(\int\nabla w\cdot\nabla v=0\)。取 \(v=w\)，得到 \(\nabla w=0\)；Poincaré 不等式推出 \(w=0\)。这个结论对不连通区域也成立，因为零边界约束在每个分支上排除了常数自由度。

\ProseParagraphBreak

内部方程则没有唯一性。例如任一调和函数都可加到 Poisson 方程的一个局部解上，而不改变源项。纯 Neumann 问题即使规定通量，仍允许每个连通分支上的加法常数。因而“取差再测试”最后得出梯度为零之后，还必须指出究竟由哪项边界条件或归一化排除了常数，不能直接省略最后一步。

\begin{proposition}[弱比较原理]
\label{b2:prop:15-weak-comparison}
设 \(u_1,u_2\in H_0^1(\Omega)\) 分别解 \(-\Delta u_i=f_i\)，且 \(\langle f_1-f_2,v\rangle\le0\) 对每个非负 \(v\in H_0^1(\Omega)\) 成立。则 \(u_1\le u_2\) 几乎处处。
\end{proposition}
\begin{proof}
令 \(w=u_1-u_2\)。Sobolev 截断性质保证 \(w_+=\max\{w,0\}\in H_0^1\)，且 \(\nabla w_+=\mathbf1_{\{w>0\}}\nabla w\)。以它测试差方程，得
\[\lVert\nabla w_+\rVert_2^2=\langle f_1-f_2,w_+\rangle\le0.\]
故 \(w_+=0\)，即所需结论。截断性质可先用光滑 Lipschitz 逼近正部函数、链式法则和 \(H^1\) 极限得到，再对 \(H_0^1\) 使用零边界逼近；这里没有要求解逐点可微。
\end{proof}

特别地，非负源项与零边界数据给出非负弱解，与第十三章经典最大值原理的符号一致。比较的结论是几乎处处关系；在尚未得到连续代表元之前，不应改说“每个点都非负”。

\subsection{能量估计}
\label{b2:subsec:150303}

取解自身作测试，Cauchy--Schwarz 的对偶形式给出
\[\lVert\nabla u\rVert_2^2=\langle f,u\rangle
 \le\lVert f\rVert_{H^{-1}}\lVert\nabla u\rVert_2.\]
若梯度非零可约去一因子；若为零估计显然成立。因此解映射 \(H^{-1}\to H_0^1\) 连续。对于两组数据，同一计算给出
\[\lVert\nabla(u_1-u_2)\rVert_2\le\lVert f_1-f_2\rVert_{H^{-1}}.\]
特别地，平滑数据的近似会在能量范数中产生相应解的收敛，这为从光滑近似过渡到粗糙数据提供了稳定性依据。

\ProseParagraphBreak

若 \(f\in L^2\)，进一步有 \(\lVert\nabla u\rVert_2\le C_P\lVert f\rVert_2\)、\(\lVert u\rVert_2\le C_P^2\lVert f\rVert_2\)。注意区域放大时 \(C_P\) 也会变化：在 \(r\Omega\) 上换元表明 \(C_P\) 按长度尺度增长。这与二阶方程解相对于源项多出两个长度因子的事实一致，估计中的几何依赖不能无故删除。

\begin{proposition}[有限维近似的最佳误差]
\label{b2:prop:15-galerkin-poisson}
设 \(V\subset H_0^1(\Omega)\) 是有限维子空间，\(u\) 为 \cref{b2:thm:15-poisson-existence} 的解。则存在唯一 \(u_V\in V\) 满足 \(\int\nabla u_V\cdot\nabla v=\langle f,v\rangle\) 对所有 \(v\in V\) 成立，而且
\[\lVert\nabla(u-u_V)\rVert_2=\inf_{v\in V}\lVert\nabla(u-v)\rVert_2.\]
\end{proposition}
\begin{proof}
限制到 \(V\) 后型仍有界、强制，故有限维方程唯一可解。两式相减得 \(\int\nabla(u-u_V)\cdot\nabla v=0\) 对所有 \(v\in V\) 成立。对任意 \(v\in V\)，分解 \(u-v=(u-u_V)+(u_V-v)\)，两项在梯度内积下正交，故误差平方等于两项范数平方之和。最小值恰在 \(v=u_V\) 达到。
\end{proof}

这个结论不宣称任一有限维空间都给出好的逼近；它说解离该空间有多近，近似解就有多好。若一列子空间递增且其并在 \(H_0^1\) 中稠密，最佳误差便趋零。递增性可改成对每个固定函数都能沿整个序列逼近的条件，但只有并集稠密不够，例如在良好空间之间反复插入零空间便会破坏收敛。存在性、离散方程的可解性与近似空间的逼近能力仍是三个可以分别核验的事实。

\ProseParagraphBreak

在区间 \((0,1)\) 上，这个结论可以算得很具体。取均匀节点 \(x_j=jh\)、\(h=1/N\)，令 \(V_h\) 为节点间分段仿射、两端为零的连续函数空间。基函数 \(\phi_j\) 在 \(x_j\) 取值一，在其余节点取值零。每个基函数只有相邻两段上的斜率非零，直接积分得
\[\int_0^1\phi_j'^2=\frac2h,\qquad
 \int_0^1\phi_j'\phi_{j+1}'=-\frac1h,\qquad
 \int_0^1\phi_j'\phi_k'=0\quad(|j-k|>1).\]
因此弱解的有限维近似由三对角线性方程组确定。矩阵的正定性并不需要另作行列式计算：任意系数向量对应的二次型就是分段仿射函数的梯度平方积分，零值只可能对应零函数。

\ProseParagraphBreak

例如取 \(f=1\)，精确弱解为 \(u(x)=x(1-x)/2\)。它在各小区间的导数平均等于两端函数值之差除以 \(h\)，而任意 \(v_h\in V_h\) 的导数在该段为常数。故 \(u\) 的节点线性插值 \(I_hu\) 满足 \(\int(u'- (I_hu)')v_h'=0\)，由唯一性即知 \(u_h=I_hu\)。在长度 \(h\) 的一段上，导数误差是以中点为零的线性函数，其平方积分为 \(h^3/12\)；函数误差则为 \((x-x_j)(x_{j+1}-x)/2\)。逐段积分并相加得到
\[\lVert u'-u_h'\rVert_2=\frac h{\sqrt{12}},\qquad
 \lVert u-u_h\rVert_2=\frac{h^2}{\sqrt{120}}.\]
这里第二个收敛速度比第一个快，来自二次函数的具体结构，不能仅从抽象存在定理读出。一般数据下的误差阶还需要解的正则性和近似空间的逼近估计；下一章提供其中的正则性工具。

\ProseParagraphBreak

弱解的误差也可以用残差直接测量。对任意 \(w\in H_0^1\) 定义 \(r=f+\Delta w\in H^{-1}\)，则 \(\langle r,v\rangle=\int\nabla(u-w)\cdot\nabla v\)。对单位梯度范数的测试函数取上确界，Cauchy--Schwarz 给出上界，而在误差非零时取 \(v=(u-w)/\lVert\nabla(u-w)\rVert_2\) 达到该界。因此 \(\lVert r\rVert_{H^{-1}}=\lVert\nabla(u-w)\rVert_2\)。对分段仿射近似，残差可能含节点通量跳跃，不能只计算每段内为零的二阶导数；负范数恰好保留这些分布贡献。

\subsection{非齐次边界数据}
\label{b2:subsec:150304}

给定 \(g\in H^1(\Omega)\)，令 \(u=g+w\)，其中 \(w\in H_0^1\)。方程等价于
\[\int_\Omega\nabla w\cdot\nabla v
  =\langle f,v\rangle-\int_\Omega\nabla g\cdot\nabla v\qquad(v\in H_0^1).\]
右端属于 \(H^{-1}\)，范数不超过 \(\lVert f\rVert_{H^{-1}}+\lVert\nabla g\rVert_2\)。因此齐次理论给出唯一 \(w\)，从而给出唯一满足 \(u-g\in H_0^1\) 的解；并有
\[\lVert\nabla w\rVert_2\le\lVert f\rVert_{H^{-1}}+\lVert\nabla g\rVert_2.\]
完整 \(H^1\) 估计还要加上 \(g\) 的 \(L^2\) 范数。

\ProseParagraphBreak

若另一个延拓 \(\widetilde g\) 与 \(g\) 相差 \(H_0^1\) 函数，两者给出的仿射可行集完全相同；由唯一性得到同一个 \(u\)。所以解依赖边界类，而不依赖人为选择的延拓。若要比较两组边界数据，先选择它们差的一个延拓，再把上面的稳定性估计应用于差方程；优化延拓的范数便自然引出商空间形式的迹范数。

\ProseParagraphBreak

例如在单位球上给定仿射边界数据 \(g(x)=a\cdot x+b\)，且 \(f=0\)。由于 \(\int a\cdot\nabla v=0\) 对 \(v\in H_0^1\) 成立，齐次修正为零，故 \(u=g\)。这里无需使用球的 Poisson 核；而对一般边界数据，弱方法仍给出存在性，之后才可能借第十三章的表示公式识别它。不同构造的解一旦处于同一能量边界类，就由唯一性自动一致。

% !TeX root = ../../Book2.tex
\section{一般二阶椭圆算子}
\label{b2:sec:1504}

Poisson 方程的证明只用到梯度型的有界性与正性，因而可以推广到系数随位置变化的算子。困难主要来自两处：低阶项可能破坏正性，非对称主部使极小化不再适用。本节先得到移位后的可解性，再在选读小节中借紧性解释原问题的有限维障碍。

\subsection{散度型算子}
\label{b2:subsec:150401}

设 \(\Omega\subset\mathbb R^d\) 有界，考虑
\[Lu=-\operatorname{div}(A\nabla u)+b\cdot\nabla u+cu,\]
其中矩阵 \(A\)、向量 \(b\) 和实函数 \(c\) 均有界可测。其 \term[sandu-xing-suanzi]{散度型算子}[divergence-form operator]主部用通量 \(A\nabla u\) 定义，不必对系数求经典导数。对应型为
\[a(u,v)=\int_\Omega A\nabla u\cdot\nabla v
       +\int_\Omega(b\cdot\nabla u)v+\int_\Omega cuv.
\]
弱 Dirichlet 问题是在 \(H_0^1(\Omega)\) 中求 \(u\)，使 \(a(u,v)=\langle f,v\rangle\) 对所有 \(v\in H_0^1\) 成立。

\ProseParagraphBreak

假设存在 \(\lambda>0\)，使对几乎处处的 \(x\) 和全部实向量 \(\xi\) 有
\[\xi^{\mathsf T}A(x)\xi\ge\lambda|\xi|^2.\]
这称为一致椭圆性。由于反对称部分的二次型为零，它等价于对称部分的相同下界；但方程本身仍保留全部矩阵。若 \(A\) 不可微，不能把 \(-\partial_i(a_{ij}\partial_j u)\) 改成 \(-a_{ij}\partial_{ij}u\) 并丢掉系数导数，这两种形式是不同的方程类别。

\subsection{有界可测系数}
\label{b2:subsec:150402}

设 \(\lVert A(x)\rVert\le\Lambda\)、\(\lVert b\rVert_\infty=B\)、\(\lVert c\rVert_\infty=C_0\)。逐项应用 Hölder 与 Poincaré 不等式，得到
\[|a(u,v)|\le(\Lambda+BC_P+C_0C_P^2)\lVert\nabla u\rVert_2\lVert\nabla v\rVert_2.\]
所以有界可测系数已经足够定义连续弱型。若 \(b=0\)、\(c\ge0\)，椭圆性还给出 \(a(u,u)\ge\lambda\lVert\nabla u\rVert_2^2\)。由 Lax--Milgram，对每个 \(f\in H^{-1}\) 唯一可解，且
\[\lVert\nabla u\rVert_2\le\lambda^{-1}\lVert f\rVert_{H^{-1}}.\]
这些结论没有要求矩阵对称。

\ProseParagraphBreak

系数粗糙时，解的导数可能在界面跳跃，\cref{b2:prop:15-interface-flux} 已给出最简单的例子。这说明存在性定理与下一章的 \(H^2\) 正则性定理必须使用不同假设：前者只要系数有界可测，后者的差商证明需要系数具有 Lipschitz 控制。不能因为方程是“椭圆的”，就把所有正则性结论一并附加上去。

\ProseParagraphBreak

对于非齐次边界数据，仍写 \(u=g+w\)。右端变成 \(f-a(g,\cdot)\)；由同样的积分估计，它是 \(H_0^1\) 上的连续泛函。因此齐次问题的结论立即推广。需要注意，估计 \(a(g,v)\) 时 \(g\) 不属于 \(H_0^1\)，不能对 \(g\) 使用零边界 Poincaré 不等式，而应保留其完整 \(H^1\) 范数。

\subsection{Gårding 不等式}
\label{b2:subsec:150403}

\begin{proposition}[Gårding 不等式]
\label{b2:prop:15-garding}
设 \(\Omega\) 有界，\(A,b,c\) 如上有界可测，且 \(A\) 以常数 \(\lambda>0\) 一致椭圆。令
\[\beta=\frac{\lVert b\rVert_\infty^2}{2\lambda}+\lVert c_-\rVert_\infty,
 \qquad c_- =\max\{-c,0\}.\]
则对所有 \(u\in H_0^1(\Omega)\)，
\[a(u,u)\ge\frac\lambda2\lVert\nabla u\rVert_2^2-\beta\lVert u\rVert_2^2.\]
因而 \(a_\mu(u,v)=a(u,v)+\mu\int_\Omega uv\) 在每个 \(\mu\ge\beta\) 时强制；方程 \((L+\mu I)u=f\) 对每个 \(f\in H^{-1}\) 唯一弱可解。
\end{proposition}
\begin{proof}
主部下界是 \(\lambda\lVert\nabla u\rVert_2^2\)。一阶项由 Young 不等式满足
\[\left|\int(b\cdot\nabla u)u\right|
 \le B\lVert\nabla u\rVert_2\lVert u\rVert_2
 \le\frac\lambda2\lVert\nabla u\rVert_2^2+\frac{B^2}{2\lambda}\lVert u\rVert_2^2.\]
零阶项下界为 \(-\lVert c_-\rVert_\infty\lVert u\rVert_2^2\)。相加即得不等式；加上 \(\mu\lVert u\rVert_2^2\) 后，强制常数至少为 \(\lambda/2\)。有界性已证明，故应用 \cref{b2:thm:15-lax-milgram} 得到最后结论。
\end{proof}

上述估计称为 \term[garding-budengshi]{Gårding 不等式}[Gårding inequality]。它不是原型的强制性，而是允许减去一个较低阶的 \(L^2\) 项。移位恢复正性，却也改变了方程：解 \((L+\mu I)u=f\) 并不能直接宣称已经解了 \(Lu=f\)。去掉移位需要额外论证，下一小节用紧嵌入完成这一步。

\ProseParagraphBreak

有时漂移项可以更精确处理。若 \(b\) 足够光滑，且 \(u\in H_0^1\)，则密度和分部积分给出 \(\int(b\cdot\nabla u)u=-\tfrac12\int(\operatorname{div}b)u^2\)。若 \(c-\tfrac12\operatorname{div}b\ge0\)，原型已强制，无须移位。对于仅有界可测的 \(b\)，散度未必是可比较的函数，此时前面的 Young 估计才是无需额外正则性的可靠结论。

\OptionalSubsection{Fredholm 择一定理}{b2:subsec:150404}

本小节只使用第一卷已经完整证明的紧算子 Fredholm 理论。应用的紧性来自第一卷推论 25.2.5：任意有界开集上的 \(H_0^1(\Omega)\hookrightarrow L^2(\Omega)\) 紧嵌入。它不需要 Lipschitz 边界，因为零延拓与平移紧性已经控制了边界附近的质量。

\begin{theorem}[椭圆边值问题的有限维障碍]
\label{b2:thm:15-fredholm-elliptic}
在 \cref{b2:prop:15-garding} 的假设下，令
\[
 \begin{aligned}
 N&=\{u\in H_0^1:a(u,v)=0\ \text{对所有 }v\in H_0^1\},\\
 N^*&=\{w\in H_0^1:a(v,w)=0\ \text{对所有 }v\in H_0^1\}.
 \end{aligned}
\]
则两空间有限维且维数相等。对 \(f\in H^{-1}\)，弱问题 \(a(u,v)=\langle f,v\rangle\) 有解，当且仅当
\[\langle f,w\rangle=0\qquad(w\in N^*).\]
有解时全部解为 \(u_0+N\)。特别地，齐次问题只有零解，当且仅当任意右端都有唯一解。
\end{theorem}
\begin{proof}
记 \(H=H_0^1\)，以梯度为内积。分别用算子 \(T,T_\mu,M:H\to H\) 表示型 \(a,a_\mu\) 及 \(\int uv\)。选 \(\mu\ge\beta\)，Lax--Milgram 使 \(T_\mu\) 成为有界可逆算子。又 \(T=T_\mu-\mu M=T_\mu(I-K)\)，其中 \(K=\mu T_\mu^{-1}M\)。

为证 \(M\) 紧，取 \(H\) 中有界列 \((u_n)\)。紧嵌入给出在 \(L^2\) 中收敛的子列；且
\[\lVert M(u_n-u_k)\rVert_H
 =\sup_{\lVert v\rVert_H\le1}\left|\int(u_n-u_k)v\right|
 \le C_P\lVert u_n-u_k\rVert_2,\]
所以像子列在 \(H\) 中收敛。由此 \(K\) 紧。第一卷定理 C.4.4 说明 \(I-K\) 有闭值域、有限维核与余核，二者维数相等；左乘可逆算子后，\(T\) 保留这些性质。

在 Hilbert 空间中，闭值域满足 \(\operatorname{ran}T=(\ker T^*)^\perp\)。而 \(\ker T=N\)、\(\ker T^*=N^*\)，因为 \(\langle Tv,w\rangle_H=a(v,w)\)。将右端 \(f\) 用 Riesz 向量表示，可解性即等价于它与 \(N^*\) 正交，也就是所列对偶配对为零。两个解之差恰在 \(N\) 内，得到所有结论。
\end{proof}

对称型的 \(N^*\) 就是 \(N\)；非对称型必须检查伴随核，不能用原核替代。例如有限维非对称矩阵的值域正交于转置矩阵的核，而未必正交于自身的核，这里的区别完全相同。定理中的相容条件只有有限多个，可以对 \(N^*\) 的一组基逐个检验。

\ProseParagraphBreak

在 \((0,\pi)\) 上取 \(L=-\frac{\mathrm d^2}{\mathrm dx^2}-1\) 和零端点条件。齐次核为 \(\operatorname{span}\{\sin x\}\)，故 \(Lu=f\) 可解恰当且仅当 \(\int_0^\pi f(x)\sin x\,\mathrm dx=0\)。源项 \(f=\sin x\) 因此不可解，而 \(f=\sin2x\) 有解 \(u=\frac13\sin2x+c\sin x\)。这个例子同时显示：一致椭圆主部并不排除共振；移位后的唯一可解性与原问题的有限维障碍可以同时成立。

% 本章小结（不计入编号）
\begin{chaptersummary}
弱形式减少了未知函数所需的导数阶数，却没有丢掉方程：紧支集测试辨认内部散度，测试空间和边界类决定所解的边值问题，界面项记录通量平衡。负阶空间把函数源项与散度型源项放在同一连续对偶中。Lax--Milgram 用有界性构造算子，以强制性控制逆，再排除值域的正交补，得到无需对称性的存在与稳定性。

\ProseParagraphBreak

一般椭圆主部只需有界可测系数即可定义能量弱解；低阶项可能使原型失去强制性，但加上足够大的零阶移位可恢复它。去掉移位时，紧嵌入将问题化为恒等算子的紧扰动，可解性由伴随核的有限个条件决定。整个论证首先产生 \(H^1\) 解；要得到 \(H^2\)、连续梯度乃至经典解，下一章还须逐项增加系数、数据和边界的正则性。
\end{chaptersummary}


\BookExercises

各题均为本书自拟。前几题核验空间和数据，随后比较近似、系数扰动与共振；最后一题使用选读小节的相容条件。

% 自拟：一维点源的精确对偶范数。
\begin{exercise}\label{b2:ex:15-point-source}
在 \((0,1)\) 上，以梯度范数定义 \(H^{-1}\)。对 \(0<s<1\)，求点质量 \(\delta_s\) 的精确范数，并说明哪个函数达到对偶上确界。
\end{exercise}
\begin{solution}
令 \(G_s(x)=(1-s)x\)（\(x\le s\)）、\(G_s(x)=s(1-x)\)（\(x\ge s\)）。正文已验证 \(v(s)=\int G_s'v'\)。因此 \(|v(s)|\le\lVert G_s'\rVert_2\lVert v'\rVert_2\)，且取 \(v=G_s/\lVert G_s'\rVert_2\) 达到上界。直接计算 \(\lVert G_s'\rVert_2^2=s(1-s)^2+(1-s)s^2=s(1-s)\)，所以范数为 \(\sqrt{s(1-s)}\)。
\end{solution}

% 自拟：高维点质量的尺度障碍。
\begin{exercise}\label{b2:ex:15-high-dimensional-point}
设 \(B(0,1)\subset\Omega\subset\mathbb R^d\) 有界。证明当 \(d\ge2\) 时点质量 \(\delta_0\notin H^{-1}(\Omega)\)。二维可使用在 \(r\le\varepsilon\) 等于 \(1\)、在 \(\varepsilon<r<1\) 等于 \(\log r/\log\varepsilon\)、在 \(r\ge1\) 等于零的径向截断。
\end{exercise}
\begin{solution}
若点质量属于该对偶，应有 \(|\varphi(0)|\le C\lVert\nabla\varphi\rVert_2\) 对所有光滑紧支集函数成立。\(d>2\) 时取 \(\varphi_\varepsilon(x)=\varphi(x/\varepsilon)\)，其中 \(\varphi(0)=1\)、支集在单位球；梯度范数平方为 \(\varepsilon^{d-2}\lVert\nabla\varphi\rVert_2^2\to0\)，矛盾。二维的题设截断 \(w_\varepsilon\) 有 \(w_\varepsilon(0)=1\)，且
\[\int|\nabla w_\varepsilon|^2=\frac{2\pi}{|\log\varepsilon|^2}\int_\varepsilon^1\frac{\mathrm dr}{r}=\frac{2\pi}{|\log\varepsilon|}\to0.\]
可先把外半径换成固定小于 \(1\) 的数，再用半径小于 \(\varepsilon/2\) 的磨光核平滑。中心仍等于 \(1\)，支集仍紧含于单位球，梯度范数不增；故得到同样矛盾。二维必须使用对数截断，简单缩放的梯度范数恰好不变。
\end{solution}

% 自拟：分段材料的显式有效导热。
\begin{exercise}\label{b2:ex:15-layered-material}
令 \(a=1\) 于 \((0,1/2)\)、\(a=4\) 于 \((1/2,1)\)。求 \(-(au')'=0\)、\(u(0)=0,u(1)=1\) 的弱解，计算左右斜率及总能量，并判断它是否属于 \(H^2(0,1)\)。
\end{exercise}
\begin{solution}
通量为常数 \(q\)，且 \(1=\int_0^1u'=q(1/2+1/8)=5q/8\)，所以 \(q=8/5\)。左斜率为 \(8/5\)，右斜率为 \(2/5\)，解是通过 \((0,0),(1/2,4/5),(1,1)\) 的分段仿射函数。能量 \(\tfrac12\int a|u'|^2=\tfrac12q\int u'=4/5\)。一阶导数在中点有非零跳跃，二阶分布导数含点质量，故不属于 \(L^2\)，解不在 \(H^2\)。
\end{solution}

% 自拟：一般强制型的Galerkin误差而非对称投影。
\begin{exercise}\label{b2:ex:15-cea}
设 \(a\) 满足 Lax--Milgram 条件，常数为 \(M,\alpha\)，\(V\subset H\) 有限维。若 \(u\in H\)、\(u_V\in V\) 分别解全空间与限制测试方程，证明
\[\lVert u-u_V\rVert_H\le\frac M\alpha\inf_{v\in V}\lVert u-v\rVert_H.\]
说明非对称情形为何不能直接引用正交投影的勾股恒等式。
\end{exercise}
\begin{solution}
限制型仍强制，故 \(u_V\) 存在唯一。误差 \(e=u-u_V\) 满足 \(a(e,w)=0\) 对所有 \(w\in V\)。于是对 \(v\in V\)，
\[\alpha\lVert e\rVert^2\le a(e,e)=a(e,u-v)+a(e,v-u_V)=a(e,u-v)
 \le M\lVert e\rVert\lVert u-v\rVert.\]
误差为零时显然，否则约去并取下确界即可。非对称型一般不是内积，测试正交并不等于 Hilbert 内积正交，故正文 Poisson 问题的勾股公式不能直接套用。
\end{solution}

% 自拟：系数在L无穷范数下的稳定性。
\begin{exercise}\label{b2:ex:15-coefficient-stability}
设 \(A_n,A\) 有共同椭圆下界 \(\lambda>0\)，且 \(\lVert A_n-A\rVert_\infty\to0\)。同一 \(f\in H^{-1}\) 对应的零边界解为 \(u_n,u\)。证明 \(u_n\to u\) 于 \(H_0^1\)，并给出定量估计。
\end{exercise}
\begin{solution}
以 \(A_n\) 为左端，两方程相减得 \(\int A_n\nabla(u_n-u)\cdot\nabla v=\int(A-A_n)\nabla u\cdot\nabla v\)。取 \(v=u_n-u\)，利用椭圆性与 Hölder 不等式，得到
\[\lVert\nabla(u_n-u)\rVert_2\le\lambda^{-1}\lVert A_n-A\rVert_\infty\lVert\nabla u\rVert_2
 \le\lambda^{-2}\lVert A_n-A\rVert_\infty\lVert f\rVert_{H^{-1}}.\]
右端趋零，Poincaré 不等式再给出完整 \(H^1\) 收敛。
\end{solution}

% 自拟：Neumann归一化与相容性。
\begin{exercise}\label{b2:ex:15-neumann}
设 \(\Omega\) 为连通有界 Lipschitz 区域，\(f\in L^2(\Omega)\)。使用零均值 Poincaré 不等式，证明 \(\int\nabla u\cdot\nabla v=\int fv\) 对所有 \(v\in H^1(\Omega)\) 可解，当且仅当 \(\int f=0\)。证明指定 \(\int u=0\) 后解唯一。
\end{exercise}
\begin{solution}
取 \(v=1\) 得必要条件。反之在闭空间 \(H=\{v\in H^1:\int v=0\}\) 上，梯度范数因零均值 Poincaré 不等式而等价于 \(H^1\) 范数，故 Lax--Milgram 给出唯一零均值解。对一般测试函数减去平均值即可落入 \(H\)，而右端常数部分因 \(\int f=0\) 消失，故方程对全部 \(H^1\) 成立。任意两个解之差梯度为零，连通性使其为常数，零均值再使该常数为零。
\end{solution}

% 自拟：漂移项的反对称结构。
\begin{exercise}\label{b2:ex:15-divergence-free-drift}
设 \(b\in C^1(\overline\Omega;\mathbb R^d)\)、\(\operatorname{div}b=0\)，\(A\) 有界且一致椭圆。证明零边界问题 \(-\operatorname{div}(A\nabla u)+b\cdot\nabla u=f\) 对任意 \(f\in H^{-1}\) 唯一可解，其能量估计不含 \(\lVert b\rVert_\infty\)。
\end{exercise}
\begin{solution}
对紧支集光滑 \(u\)，有 \(\int(b\cdot\nabla u)u=\tfrac12\int b\cdot\nabla(u^2)=0\)。双线性连续性和 \(H_0^1\) 密度把该等式推广到全部 \(H_0^1\)。所以型虽可能不对称，却满足 \(a(u,u)\ge\lambda\lVert\nabla u\rVert_2^2\)。有界性仍含 \(\lVert b\rVert_\infty\)，但强制常数不含它。Lax--Milgram 给出 \(\lVert\nabla u\rVert_2\le\lambda^{-1}\lVert f\rVert_{H^{-1}}\) 及唯一性。
\end{solution}

% 自拟：不同共振频率的相容条件与全体解。
\begin{optionalexercise}\label{b2:ex:15-resonance}
在 \((0,\pi)\) 上求解零端点问题 \(-u''-4u=\sin x+q\sin2x\)。确定哪些 \(q\in\mathbb R\) 允许弱解，并在可解时写出全部解。
\end{optionalexercise}
\begin{solution}
型对称，齐次核是 \(\operatorname{span}\{\sin2x\}\)。与核配对得相容条件 \(q\int_0^\pi\sin^22x=0\)，即 \(q=0\)。此时特解为 \(-\tfrac13\sin x\)，故全部解为 \(-\tfrac13\sin x+c\sin2x\)。核的刻画可由一维分布方程先推出 \(u''=-4u\in L^2\)，再由常系数常微分方程及端点条件得到，亦可直接使用正弦展开。因此没有遗漏较粗糙的能量解。
\end{solution}

% !TeX root = ../../Book2.tex
\chapter{椭圆正则性}
\label{b2:ch:16}
\BookChapterNumber{b2:ch:16}
% 本章摘要（不计入编号）
\begin{chapterabstract}
本章研究能量弱解何时具有更高正则性。差商方法在不预设二阶导数的前提下，建立 Lipschitz 系数方程的内部 \(H^2\) 估计；逐项控制求导误差后，得到高阶 Sobolev 正则性和光滑数据的内部光滑性。边界部分通过平坦化、切向差商和方程恢复法向导数，证明 \(C^{1,1}\) 区域上的 Dirichlet 二阶估计。最后仅用有界可测系数建立 Caccioppoli 不等式及截断迭代的局部有界性，并说明 Harnack 比较的完整线性证明如何在附录 H 补齐，随后推出 Hölder 连续性。
\end{chapterabstract}

% 本章目标（不计入编号）
\begin{learninggoals}
\begin{itemize}
\item\label{b2:item:16-goal-difference}使用离散分部积分和一致差商界构造二阶弱导数，指出系数导数的作用。
\item\label{b2:item:16-goal-bootstrap}逐步核验微分后右端的空间，结合 Sobolev 嵌入判断经典正则性。
\item\label{b2:item:16-goal-boundary}理解边界平坦化、切向估计及纯法向恢复，并用凹角反例检验全局结论。
\item\label{b2:item:16-goal-low}在可测系数下证明局部能量界和有界性，区分本章已证结论与附录 H 补齐的 Harnack 比较证明。
\end{itemize}
\end{learninggoals}

上一章的弱解可以没有经典二阶导数，但方程并非只是一条约束：它还把不同方向和不同尺度的变化联系起来。若系数变化受控，差商误差可以吸收，解便得到新的弱导数；若系数只是可测，导数提升可能失败，标量解的高值集合却仍受能量限制。本章分别建立这两类估计，避免把“椭圆正则性”当作无条件的光滑性口号。

\ProseParagraphBreak

本章普通主线承接第十五章的测试等式，使用第一卷第二十四章的 Sobolev 与 Poincaré 不等式。系统参考可读 Lawrence C. Evans 的《Partial Differential Equations》\cite{Evans2010Partial}以及 David Gilbarg 与 Neil S. Trudinger 的《Elliptic Partial Differential Equations of Second Order》\cite{Gilbarg2001Elliptic}。下面对差商、边界 \(H^2\) 与局部有界性给出完整证明；末节先给可测系数 Harnack 比较的证明概要，完整的线性比较证明及其所需指数尾见附录 H、C；不把此处略去的步骤隐藏在“标准估计”中。

% !TeX root = ../../Book2.tex
\section{差商方法}
\label{b2:sec:1601}

能量解只有一阶弱导数，不能直接对方程再求导。差商以两个平移值之差代替导数，先在已有的函数空间中操作，再由一致估计恢复二阶弱导数。这样既避免预设待证正则性，也清楚显示系数的变化会在估计中产生什么误差。

\subsection{差商}
\label{b2:subsec:160101}

设整数 \(d\geq1\)，\(\Omega\subset\mathbb R^d\) 为开集，\(u\in L^1_{\mathrm{loc}}(\Omega;\mathbb K)\)，其中 \(\mathbb K=\mathbb R\) 或 \(\mathbb C\)，\(k\in\{1,\ldots,d\}\)、\(h\in\mathbb R\setminus\{0\}\)。固定坐标方向 \(e_k\)，定义 \term[chashang]{差商}[difference quotient]
\[D_h^ku(x)=\frac{u(x+he_k)-u(x)}h\qquad(h\ne0).\]
只在所需平移仍留在定义域内时使用这一表达式。若 \(f,g\in C_c^\infty(\Omega)\)，并且它们支集的所需 \(h\) 平移都仍在 \(\Omega\) 内，换元给出离散分部积分；该式也适用于满足这些支撑条件的 \(L^2\) 函数
\[\int fD_{-h}^kg=-\int(D_h^kf)g.\]
其中 \(D_{-h}^kg=(g(x-he_k)-g(x))/(-h)\)；负号和负步长必须同时处理。乘积公式是
\[D_h^k(fg)=f(x+he_k)D_h^kg+(D_h^kf)g(x).\]

\Needspace{6\baselineskip}
\begin{lemma}[差商与弱导数]
\label{b2:lem:16-difference-quotient}
设整数 \(d\geq1\)，\(U\Subset V\subset\mathbb R^d\) 为非空开集，标量域 \(\mathbb K=\mathbb R\) 或 \(\mathbb C\)，\(k\in\{1,\ldots,d\}\)，\(h\in\mathbb R\setminus\{0\}\)，\(D_h^ku(x)=[u(x+he_k)-u(x)]/h\)，\(|h|<\operatorname{dist}(U,\partial V)\)（\(V=\mathbb R^d\) 时该距离取 \(+\infty\)）。若 \(u\in H^1(V;\mathbb K)\)，则
\[\lVert D_h^ku\rVert_{L^2(U)}\le\lVert\partial_ku\rVert_{L^2(V)}.\]
反之，若 \(u\in L^2(V;\mathbb K)\)，且存在 \(0\leq C<\infty\)、\(h_0>0\)，使 \(\lVert D_h^ku\rVert_{L^2(U)}\le C\) 对所有 \(0<|h|<h_0\) 且 \(|h|<\operatorname{dist}(U,\partial V)\) 成立，则 \(u\) 在 \(U\) 中有弱导数 \(\partial_ku\in L^2(U;\mathbb K)\)，其范数不超过 \(C\)。
\end{lemma}
\begin{proof}
对光滑 \(u\)，有 \(D_h^ku(x)=\int_0^1\partial_ku(x+the_k)\,\mathrm dt\)。Jensen 不等式、积分换序和各平移集合包含于 \(V\) 给出所列界。一般 \(H^1\) 函数先在包住所有相关线段的内部区域磨光，再以 \(H^1\) 收敛取极限，结论保持。

证明反向时，取 \(h_n\to0\)。\(L^2(U)\) 中的有界差商有弱收敛子列，记极限为 \(g\)。对 \(\varphi\in C_c^\infty(U)\)，离散分部积分和测试函数差商的一致收敛给出
\[\int_Ug\varphi=\lim_n\int_U(D_{h_n}^ku)\varphi
 =-\lim_n\int_VuD_{-h_n}^k\varphi=-\int_Uu\partial_k\varphi.\]
故 \(g\) 正是弱导数，弱下半连续性给出范数界。弱导数的唯一性说明所获对象不依赖抽取的子列。
\end{proof}

引理中的第一项是已有导数给出的上界，第二项则把统一差商界变成导数存在性；正则性证明需要把这两个方向配合使用。仅对某一个固定步长估计当然不能推出可微性。还应保留内部集合 \(U\Subset V\)：在边界附近任意平移会跑出区域，边界估计因此需要新的处理。

\subsection{内部 \texorpdfstring{\(H^2\)}{H2} 估计}
\label{b2:subsec:160102}

\Needspace{6\baselineskip}
\begin{theorem}[内部二阶估计]
\label{b2:thm:16-interior-h2}
设整数 \(d\geq1\)，\(\Omega\subset\mathbb R^d\) 为非空开集，\(r>0\)，\(B_{2r}=B(x_0,2r)\Subset\Omega\)，实矩阵场 \(A\in W^{1,\infty}(B_{2r};\mathbb R^{d\times d})\) 满足 \(\lVert A\rVert_\infty\le\Lambda<\infty\)，且对几乎处处 \(x\in B_{2r}\)、全部 \(\xi\in\mathbb R^d\) 有 \(\xi^{\mathsf T}A(x)\xi\ge\lambda|\xi|^2\)，其中 \(\lambda>0\)，矩阵取欧氏算子范数。若实值 \(u\in H^1(B_{2r};\mathbb R)\)、\(f\in L^2(B_{2r};\mathbb R)\) 且
\[\int_{B_{2r}}A\nabla u\cdot\nabla\varphi=\int_{B_{2r}}f\varphi
 \qquad(\varphi\in C_c^\infty(B_{2r};\mathbb R)),\]
则 \(u\in H^2(B_r)\)，并有
\[\lVert D^2u\rVert_{L^2(B_r)}
 \le C\left(\lVert f\rVert_{L^2(B_{2r})}
 +(r^{-1}+\lVert\nabla A\rVert_\infty)\cdot\lVert\nabla u\rVert_{L^2(B_{2r})}\right),\]
其中 \(C\) 只依赖 \(d,\lambda,\Lambda\)。
\end{theorem}
\begin{proof}
取光滑截断 \(0\le\eta\le1\)，使其在 \(B_r\) 上为 \(1\)，支集包含于 \(B_{3r/2}\)，且 \(|\nabla\eta|\le C_d/r\)。固定方向 \(k\)，取 \(|h|<r/8\)，记 \(w=D_h^ku\)。对这个固定步长，\(w\) 属于较小球上的 \(H^1\)，所以 \(v=-D_{-h}^k(\eta^2w)\) 属于 \(H_0^1(B_{2r})\)。型与右端均连续，密度保证可以用 \(v\) 测试。

将差商移到通量，得到
\[
 \int D_h^k(A\nabla u)\cdot\nabla(\eta^2w)
 =-\int fD_{-h}^k(\eta^2w).
\]
乘积公式给出 \(D_h^k(A\nabla u)=A(x+he_k)\nabla w+(D_h^kA)\nabla u\)。其中主项为 \(\int\eta^2A(x+he_k)\nabla w\cdot\nabla w\)，下界是 \(\lambda X^2\)，记 \(X=\lVert\eta\nabla w\rVert_2\)。设 \(U=\lVert\nabla u\rVert_{L^2(B_{2r})}\)、\(F=\lVert f\rVert_{L^2(B_{2r})}\)、\(K=\lVert\nabla A\rVert_\infty\)。引理给出 \(\lVert w\rVert_{L^2(B_{7r/4})}\le U\)，且 Lipschitz 控制给出 \(\lVert D_h^kA\rVert_\infty\le K\)。

展开 \(\nabla(\eta^2w)=\eta^2\nabla w+2\eta w\nabla\eta\)。其余主部交叉项、系数差商项及右端分别被
\[C\Lambda r^{-1}UX,\qquad KUX+CKr^{-1}U^2,
 \qquad CF(X+r^{-1}U)\]
控制。最后一界再次使用差商引理于紧支集函数 \(\eta^2w\)，即 \(\lVert D_{-h}^k(\eta^2w)\rVert_2\le\lVert\partial_k(\eta^2w)\rVert_2\)。逐项用 Young 不等式吸收不超过 \(\lambda X^2/2\) 的部分，余项得到
\[X^2\le C\bigl(F^2+(r^{-1}+K)^2U^2\bigr),\]
常数与 \(h\) 无关。因为 \(\eta=1\) 于 \(B_r\)，各 \(D_h^k\partial_j u\) 在该球的 \(L^2\) 范数一致有界。应用引理的反向得到全部 \(\partial_k\partial_j u\in L^2(B_r)\)，并对有限个指标求和即得估计。
\end{proof}

这个证明没有先假定二阶导数存在。测试函数中的差商仍只要求 \(u\in H^1\)；真正的新导数是在获得与步长无关的界之后才出现。系数的 Lipschitz 范数也有明确来源：它控制 \(D_hA\)。若系数只有本性上界，这个误差可能按 \(|h|^{-1}\) 增大，证明就不能闭合。

这里的矩阵系数取实值，解和右端的复值版本可直接分开处理：若 \(u=u_1+\mathrm i u_2\)、\(f=f_1+\mathrm i f_2\)，同一实矩阵 \(A\) 使两个实方程分别成立。各用一次定理，再把两个平方估计相加，就得到同一形式、仍只依赖 \(d,\lambda,\Lambda\) 的复值估计。下节的高阶提升也可逐次按这个办法处理复值数据；一般复矩阵的椭圆条件与配对须另作规定。

\subsection{局部正则性}
\label{b2:subsec:160103}

设整数 \(d\geq1\)，\(\Omega\subset\mathbb R^d\) 为非空开集，\(A\in W_{\mathrm{loc}}^{1,\infty}(\Omega;\mathbb R^{d\times d})\)，实值 \(f\in L_{\mathrm{loc}}^2(\Omega)\)、\(u\in H_{\mathrm{loc}}^1(\Omega)\) 满足 \(-\operatorname{div}(A\nabla u)=f\) 的局部测试恒等式。给定非空开集 \(U\Subset V\Subset\Omega\)，假设对几乎处处 \(x\in V\)、全部 \(\xi\in\mathbb R^d\) 有 \(\xi^{\mathsf T}A(x)\xi\geq\lambda_V|\xi|^2\)，其中 \(\lambda_V>0\)。用等半径内部球覆盖 \(\overline U\)，使加倍球均在 \(V\) 内、重叠次数只由 \(d\) 控制，再将球上估计平方求和，得到
\[\lVert u\rVert_{H^2(U)}\le C_{U,V,A}\bigl(\lVert f\rVert_{L^2(V)}+\lVert u\rVert_{H^1(V)}\bigr).\]
常数可只按 \(d\)、\(\operatorname{dist}(U,\partial V)^{-1}\)、\(\lambda_V\) 及 \(\|A\|_{W^{1,\infty}(V)}\) 选取。若每个紧含子域都有这样的椭圆下界，则局部 Lipschitz 系数和局部平方可积右端把能量解提高到 \(H_{\mathrm{loc}}^2\)。常数在区域逼近边界时可能发散，不能由此直接得到全域 \(H^2\) 估计。

\ProseParagraphBreak

若方程还含 \(b\cdot\nabla u+cu\)，且低阶系数有界，则把它移到右端：\(\widetilde f=f-b\cdot\nabla u-cu\in L_{\mathrm{loc}}^2\)。内部二阶估计仍适用，右端多出已有的 \(H^1\) 范数。这一步不需要对 \(b,c\) 求导；只有继续提高阶数时才会用到它们更高的正则性。

\ProseParagraphBreak

一个对照例是前章的分段导热系数：弱解通量连续而梯度跳跃，故二阶分布导数含点质量，不在 \(L^2\)。它说明系数正则性要求具有真实数学内容。另一方面，对常系数 Laplace 方程，本定理不要求 \(f\) 连续，只要求 \(f\in L^2\)；所得结论是二阶弱导数平方可积，而不是二阶经典导数连续。正则性结论应始终与数据所在的空间一起读取。

% !TeX root = ../../Book2.tex
\section{自举法}
\label{b2:sec:1602}

得到二阶弱导数后，方程可以合法地再作分布微分。但每次微分都会同时作用于系数，产生新的右端。自举的关键不是一句“反复求导”，而是证明这些新项具有足够的可积性，使上一节的估计能够再次使用。

\subsection{系数正则性}
\label{b2:subsec:160201}

以 \(-\operatorname{div}(A\nabla u)=f\) 为例，形式求 \(x_k\) 导数得到
\[
 -\operatorname{div}(A\nabla\partial_ku)
 =\partial_kf+\operatorname{div}((\partial_kA)\nabla u).
\]
若只有 \(A\in W^{1,\infty}\)、\(f\in L^2\)，右端一般只属于 \(H^{-1}\)，不能再次套用要求 \(L^2\) 右端的二阶定理。若进一步 \(A\in W^{2,\infty}\)、\(f\in H^1\)，并且已经知道 \(u\in H^2\)，则 \((\partial_kA)\nabla u\in H^1\)，其散度属于 \(L^2\)，因而可以把 \(\partial_ku\) 提高到 \(H^2\)。

\ProseParagraphBreak

系数的这一级额外导数不可随意省略。例如在一维取正的 Lipschitz 系数 \(a(x)=1+|x|\)，令 \(u'(x)=1/a(x)\)。则 \(-(au')'=0\)，\(u\in H_{\mathrm{loc}}^2\)，但 \(u''\) 在零点有跳跃，故 \(u\notin H_{\mathrm{loc}}^3\)。右端已光滑，解仍受到系数不光滑的限制。这个例子精确对应上一行中对 \(\partial_kA\) 再求导的需要。

\ProseParagraphBreak

本节所列 \(W^{m+1,\infty}\) 条件是便于统一证明的一组充分条件，并非每个具体方程的最弱假设。常系数的导数项全部消失；若某些系数仅沿特定方向变化，也可能先在其他方向取得更高正则性。使用定理时可以利用这些特殊结构，但需要重做相应的乘积估计，不能把一般结论的假设凭经验删除。

\subsection{高阶 Sobolev 正则性}
\label{b2:subsec:160202}

\begin{theorem}[内部高阶正则性]
\label{b2:thm:16-higher-regularity}
设 \(m\) 为非负整数，\(A\in W_{\mathrm{loc}}^{m+1,\infty}(\Omega;\mathbb R^{d\times d})\) 局部有界且一致椭圆，\(f\in H_{\mathrm{loc}}^m(\Omega)\)。若 \(u\in H_{\mathrm{loc}}^1(\Omega)\) 弱解 \(-\operatorname{div}(A\nabla u)=f\)，则 \(u\in H_{\mathrm{loc}}^{m+2}(\Omega)\)。对 \(U\Subset V\Subset\Omega\)，有
\[\lVert u\rVert_{H^{m+2}(U)}
 \le C\bigl(\lVert f\rVert_{H^m(V)}+\lVert u\rVert_{H^1(V)}\bigr),\]
其中 \(C\) 依赖区域间距、椭圆下界和 \(A\) 在 \(V\) 上的相应系数范数。
\end{theorem}
\begin{proof}
\(m=0\) 是 \cref{b2:thm:16-interior-h2} 及有限覆盖的结论。归纳假设已给出 \(u\in H_{\mathrm{loc}}^{m+1}\)。对 \(|\beta|=m\)，令 \(w=D^\beta u\in H_{\mathrm{loc}}^1\)。分布的乘积法则给出
\[
 -\operatorname{div}(A\nabla w)
 =D^\beta f+\sum_{0<\gamma\le\beta}\binom\beta\gamma
       \operatorname{div}\bigl((D^\gamma A)\nabla D^{\beta-\gamma}u\bigr).
\]
当 \(\beta=0\) 时和为空，已经处理。对和内每一项，系数 \(D^\gamma A\) 及其一阶导数有界，而 \(\nabla D^{\beta-\gamma}u\) 及其一阶导数的阶数最多为 \(m+1\)，故属于局部 \(L^2\)。乘积因此属于局部 \(H^1\)，其散度属于局部 \(L^2\)。右端第一项也由 \(f\in H^m\) 属于局部 \(L^2\)。将二阶定理应用于 \(w\)，得 \(D^\beta u\in H_{\mathrm{loc}}^2\)，即 \(u\in H_{\mathrm{loc}}^{m+2}\)。

为得到定量界，在 \(U,V\) 之间选取有限个逐层嵌套的开集。每步二阶估计的右端由前一步已经控制的 \(H^{j+1}\) 范数、\(f\) 的 \(H^j\) 范数及系数的 \(W^{j+1,\infty}\) 范数控制。依次代入归纳估计并合并有限个常数，得到所列不等式。乘积法则可以先局部磨光，再在相应 Sobolev 范数中取极限，所以没有预设经典高阶导数。
\end{proof}

对常系数 Laplace 方程，上述论证尤其直接：\(-\Delta D^\beta u=D^\beta f\)。一般变系数方程的额外项只含较低阶的解导数，正是归纳能够闭合的原因；若求导后又在右端出现尚未控制的最高阶导数，就必须先把它保留在主部，不能把它当成已知数据。

\ProseParagraphBreak

若包含低阶项且 \(b,c\) 具有相应有界导数，同样的归纳适用。对非线性方程则须额外检查复合函数的 Sobolev 正则性：从 \(u\in H^2\) 不能在所有维数自动推出 \(F(u)\in H^1\) 或 \(H^2\)。例如含高次幂的项，其导数是 \(F'(u)\nabla u\)，必须先用嵌入定理保证乘积可积。自举是一系列空间之间的合法映射，而不是只按形式导数计数。

\subsection{从弱解到经典解}
\label{b2:subsec:160203}

Sobolev 正则性最终通过嵌入转成经典正则性。第一卷定理 24.4.3 及处理正整数差的命题 24.4.4，在局部截断后给出：若整数 \(s>k+d/2\)，则 \(H_{\mathrm{loc}}^s\) 函数有 \(C^k\) 代表元；更精细地，可取得所有满足 \(k+\alpha<s-d/2\)、\(0<\alpha<1\) 的局部 \(C^{k,\alpha}\) 正则性。严格不等号不能随意替成等号。

\begin{corollary}[光滑数据的内部光滑性]
\label{b2:cor:16-smooth-interior}
设 \(A,f\in C^\infty(\Omega)\)，\(A\) 在每个紧子集附近一致椭圆。若 \(u\in H_{\mathrm{loc}}^1(\Omega)\) 弱解 \(-\operatorname{div}(A\nabla u)=f\)，则 \(u\) 有 \(C^\infty(\Omega)\) 代表元，并逐点满足方程。
\end{corollary}
\begin{proof}
固定 \(U\Subset\Omega\)。光滑系数在稍大紧集上的各阶导数有界，高阶正则性定理给出任意阶的局部 \(H^s\) 正则性。给定 \(k\)，选整数 \(s>k+d/2\)，嵌入得到 \(C^k\) 代表元；各代表元因几乎处处相同且连续而彼此一致，故得到单一 \(C^\infty\) 代表元。该代表元对应的微分表达式连续，与零分布相等，所以处处为零。
\end{proof}

若只知道 \(f\in L^2\)，在 \(d=2,3\) 时 \(H_{\mathrm{loc}}^2\) 已给出一定 Hölder 连续性，却一般不给出连续一阶导数；在 \(d\ge4\) 时甚至不能单凭这一 Sobolev 结论推出连续性。维数越高，要通过平方可积导数达到同一经典阶数，就需要越多次提升。低正则椭圆理论提供了另一条路：即使没有系数导数，也能对实标量方程控制解的振荡；这将在本章最后说明。

\ProseParagraphBreak

“弱解变成经典解”还需要明确所指区域。这里所有结论均为内部结论：边界值不光滑、区域有角点时，解可能在每个内部点都光滑，却在靠近边界处出现不可积的二阶导数。下一节将通过边界坐标和一个扇形反例说明这种差别。

% !TeX root = ../../Book2.tex
\section{边界正则性}
\label{b2:sec:1603}

内部差商可以向任意坐标方向平移，靠近边界时却可能越界。对于零 Dirichlet 条件，平坦边界允许沿切向作差商，而法向二阶导数则由方程补回。曲边界先化成平坦边界；这个坐标变换消耗的正则性，决定了边界应满足的假设。

\subsection{边界平坦化}
\label{b2:subsec:160301}

称有界区域为 \(C^{1,1}\) 区域，是指边界在有限个局部正交坐标中可写成 \(x_d=\psi(x')\)，其中 \(\psi\) 一阶可微且梯度 Lipschitz 连续，区域位于图像的一侧。局部使用
\[x=\Phi(y)=(y',y_d+\psi(y')),\qquad y_d>0\]
把半球或半盒映到区域内部。映射及其逆一阶导数有界，二阶弱导数也有界；因此它们保留局部 \(H^1\)、\(H^2\) 正则性，并给出相应范数的双向控制。

\ProseParagraphBreak

若 \(v(y)=u(\Phi(y))\)、\(\zeta(y)=\varphi(\Phi(y))\)，链式法则和换元把型写成
\[
 \int A(x)\nabla_xu\cdot\nabla_x\varphi\,\mathrm dx
 =\int\widetilde A(y)\nabla_yv\cdot\nabla_y\zeta\,\mathrm dy,
\]
其中
\[\widetilde A=|\det D\Phi|\,D\Phi^{-1}(A\circ\Phi)D\Phi^{-\mathsf T},
 \qquad\widetilde f=(f\circ\Phi)|\det D\Phi|.\]
图像平移的行列式为 \(1\)，保留一般写法便于检查其他坐标。由于 \(D\Phi^{-\mathsf T}\xi\) 的长度与 \(|\xi|\) 双向可比，\(\widetilde A\) 仍一致椭圆；若 \(A\in W^{1,\infty}\)、\(\Phi\in C^{1,1}\)，则 \(\widetilde A\in W^{1,\infty}\)。右端 \(\widetilde f\) 仍属于 \(L^2\)。这正是下一步差商证明需要的三个条件。

\ProseParagraphBreak

只有 Lipschitz 边界时，\(D\Phi\) 一般只有本性有界性，变换后的系数未必具有有界一阶导数。此时切向差商中的系数误差便不再受控。边界光滑性不是为了画图方便，而是通过坐标系数直接进入估计。

\subsection{Dirichlet 问题}
\label{b2:subsec:160302}

先考虑半球 \(B_2^+=B(0,2)\cap\{x_d>0\}\)。设 \(u\in H^1(B_2^+)\) 在平坦部分的迹为零，且弱解 \(-\operatorname{div}(A\nabla u)=f\)。测试函数可以支撑到平坦部分，但须在那里取零迹，并在球面部分附近消失。这样的测试等式由原 Dirichlet 弱形式经局部化得到。

\begin{lemma}[边界切向估计]
\label{b2:lem:16-tangential-estimate}
设 \(A\in W^{1,\infty}(B_2^+)\) 有界且一致椭圆，\(f\in L^2(B_2^+)\)，\(u\) 如上具有零平坦边界迹。则对 \(k<d\) 及 \(j\le d\)，\(\partial_k\partial_j u\in L^2(B_1^+)\)，且
\[\sum_{k<d,j\le d}\lVert\partial_k\partial_j u\rVert_{L^2(B_1^+)}
 \le C\bigl(\lVert f\rVert_{L^2(B_2^+)}+\lVert u\rVert_{H^1(B_2^+)}\bigr).\]
\end{lemma}
\begin{proof}
取在 \(B_1\) 上为 \(1\)、支集在 \(B_{3/2}\) 内的光滑截断 \(\eta\)。对切向方向 \(e_k\)，平移不改变 \(x_d\)，所以 \(D_h^ku\) 的平坦边界迹仍为零。于是 \(-D_{-h}^k(\eta^2D_h^ku)\) 是合法零迹测试函数；它在曲面边界附近为零，没有额外边界项。

在半球上重复内部估计的离散分部积分。所有平移都沿切向，故积分换元不越过平坦边界；截断又避免了球面边界。主项仍有椭圆下界，系数差商仍由 \(\lVert\nabla A\rVert_\infty\) 控制，右端仍由差商的 \(H^1\) 界控制。因此 \(\eta\nabla D_h^ku\) 的 \(L^2\) 范数独立于 \(h\) 有界。差商引理在半球内部测试并取极限，得到所列切向及混合导数及范数界。
\end{proof}

现在只缺 \(\partial_{dd}u\)。在分布意义展开方程，得
\[
 a_{dd}\partial_{dd}u
 =-f-\sum_{(i,j)\ne(d,d)}a_{ij}\partial_{ij}u
       -\sum_{i,j}(\partial_i a_{ij})\partial_j u.
\]
右端各项均属于 \(L^2(B_1^+)\)：所有非纯法向二阶项已由引理控制，系数导数有界，解的一阶导数平方可积。又 \(a_{dd}\ge\lambda>0\)，故除以它得到 \(\partial_{dd}u\in L^2\)。展开与除法均可在分布中合法进行，因为 \(a_{dd}\) 及其倒数是局部 Lipschitz 乘子。至此得到局部边界 \(H^2\) 估计。

\ProseParagraphBreak

这一步没有沿法向把 \(u\) 随意平移，也没有先把 \(u\) 作二阶可微延拓。零迹保证切向测试合法，椭圆性保证纯法向最高阶项的系数可逆，两者承担不同职责。一维时没有切向方向，最后的方程直接恢复二阶导数，恰与一维计算一致。

\subsection{边界 \texorpdfstring{\(H^2\)}{H2} 估计}
\label{b2:subsec:160303}

\begin{theorem}[全局 Dirichlet 二阶正则性]
\label{b2:thm:16-boundary-h2}
设 \(\Omega\subset\mathbb R^d\) 是有界 \(C^{1,1}\) 区域，\(A\in W^{1,\infty}(\Omega)\) 有界且一致椭圆，\(f\in L^2(\Omega)\)。若 \(u\in H_0^1(\Omega)\) 弱解 \(-\operatorname{div}(A\nabla u)=f\)，则 \(u\in H^2(\Omega)\)，且
\[\lVert u\rVert_{H^2(\Omega)}\le C\bigl(\lVert f\rVert_2+\lVert u\rVert_{H^1(\Omega)}\bigr)\le C'\lVert f\rVert_2.\]
常数依赖区域的有限边界坐标、系数上界与 Lipschitz 范数及椭圆下界。
\end{theorem}
\begin{proof}
覆盖 \(\partial\Omega\) 的有限个 \(C^{1,1}\) 坐标把方程变成上一小节的半球模型。坐标变换保留零迹、\(L^2\) 右端及 Lipschitz 椭圆系数；切向估计与恢复纯法向导数给出每个较小边界坐标内的 \(H^2\) 界。余下的紧内部区域由有限个球覆盖，应用内部二阶估计。将有限个局部积分相加，并使用坐标变换的双向范数控制，得到第一项不等式。上一章的能量估计与 Poincaré 不等式控制 \(\lVert u\rVert_{H^1}\le C\lVert f\rVert_2\)，得到第二项。
\end{proof}

对非齐次边界，若给定延拓 \(g\in H^2(\Omega)\)，令 \(w=u-g\)。右端变为 \(f+\operatorname{div}(A\nabla g)\in L^2\)，故
\[\lVert u\rVert_{H^2}\le C\bigl(\lVert f\rVert_2+\lVert g\rVert_{H^2}\bigr).\]
这里比弱存在性所需的 \(g\in H^1\) 多一阶正则性。若只有粗糙边界数据，内部解仍可能光滑，但上述全局二阶结论不再由本定理保证。

\ProseParagraphBreak

对足够光滑的区域和数据，可以在边界坐标中继续切向微分，再用方程恢复高阶法向导数。每次操作都要控制新的系数和坐标导数，因此需要同步提高边界光滑性。全局高阶定理不能仅由内部自举一句带过；本章只完整证明后续最常用的 \(H^2\) 层次。

\subsection{边界光滑性的作用}
\label{b2:subsec:160304}

取平面扇形 \(0<r<1\)、\(0<\theta<\omega\)，其中 \(\pi<\omega<2\pi\)，令 \(\alpha=\pi/\omega\in(1/2,1)\)。函数
\[v(r,\theta)=r^\alpha\sin(\alpha\theta)\]
在扇形内部调和，并在两条直边上为零。它的梯度量级为 \(r^{\alpha-1}\)，故近角点的梯度平方积分由 \(\int_0^\varepsilon r^{2\alpha-1}\,\mathrm dr\) 控制而有限；二阶导数量级为 \(r^{\alpha-2}\)，相应积分 \(\int_0^\varepsilon r^{2\alpha-3}\,\mathrm dr\) 发散。

\ProseParagraphBreak

选径向光滑截断 \(\chi\)，使它在角点附近为 \(1\)，在外圆边界附近为零，令 \(u=\chi v\)。则 \(u\in H_0^1\)，但 \(u\notin H^2\)；源项 \(f=-\Delta u\) 只支撑在离开角点的环带内，因而光滑且平方可积。局部化并在内圆 \(r=\varepsilon\) 上积分，边界余项按 \(\varepsilon^\alpha\) 消失，验证 \(-\Delta u=f\) 确实弱成立。于是即便主部是常系数、源项很好，凹角仍可破坏全局 \(H^2\) 正则性。

\ProseParagraphBreak

这个区域具有 Lipschitz 边界，却在凹角处不属于 \(C^{1,1}\)。例子不意味着所有非 \(C^{1,1}\) 区域都失败，例如某些凸多边形仍有更好的结论；它说明仅凭任意 Lipschitz 边界不能保证本节的全局二阶估计。内部估计在离角点有正距离的每个子域上仍完全成立，两个结论并不冲突。

% !TeX root = ../../Book2.tex
\section{Caccioppoli 不等式与低正则理论初步}
\label{b2:sec:1604}

有界可测系数不足以支持前面的差商估计，却仍保留椭圆性带来的能量控制。对实标量方程，截断与 Sobolev 不等式还能把积分控制逐步变成点态控制。本节完整证明局部有界性，随后明确区分已经完成的能量论证和通向 Harnack、Hölder 理论的深层步骤。

\subsection{Caccioppoli 不等式}
\label{b2:subsec:160401}

以下矩阵 \(A\) 仅假定可测，满足 \(\lVert A\rVert_\infty\le\Lambda\) 和 \(\xi^{\mathsf T}A(x)\xi\ge\lambda|\xi|^2\)。系数不需要一阶导数。

\begin{proposition}[内部能量估计]
\label{b2:prop:16-caccioppoli}
设 \(B_R\Subset\Omega\)，\(u\in H^1(B_R)\) 弱解 \(-\operatorname{div}(A\nabla u)=f\)，\(f\in L^2(B_R)\)。则对 \(0<r<R\) 和任意常数 \(c\in\mathbb R\)，
\[
 \int_{B_r}|\nabla u|^2\le
 \frac{C}{(R-r)^2}\int_{B_R}|u-c|^2
       +C(R-r)^2\int_{B_R}|f|^2,
\]
其中 \(C\) 只依赖 \(d,\lambda,\Lambda\)。
\end{proposition}
\begin{proof}
取 \(\eta\in C_c^\infty(B_R)\)，在 \(B_r\) 上为 \(1\)，满足 \(0\le\eta\le1\)、\(|\nabla\eta|\le C_d/(R-r)\)。函数 \(\eta^2(u-c)\) 是合法测试函数。展开并用椭圆性，得
\[\lambda\int\eta^2|\nabla u|^2
 \le2\Lambda\int\eta|\nabla u||u-c||\nabla\eta|
       +\int|f|\eta^2|u-c|.\]
第一项由 Young 不等式吸收一半左端，余下是 \(C\int|u-c|^2|\nabla\eta|^2\)。对第二项使用带尺度的 Young 不等式，界为 \(C(R-r)^2\int|f|^2+C(R-r)^{-2}\int|u-c|^2\)。整理并利用 \(\eta=1\) 于内球即可。
\end{proof}

这项 \term[caccioppoli-budengshi]{Caccioppoli 不等式}[Caccioppoli inequality]以较大区域内函数的大小控制较小区域内的梯度。取常数 \(c\) 可以去掉与方程无关的整体平移；齐次方程中常取 \(c\) 为球平均值，以突出振荡而不是绝对高度。估计方向与 Poincaré 不等式相反，但需要解满足方程，因此不能对任意 \(H^1\) 函数成立。

\subsection{反向估计}
\label{b2:subsec:160402}

对于齐次解，在 \(B_{2r}\) 上取 \(c=u_{B_{2r}}\)。Caccioppoli 与球上的 Sobolev--Poincaré 不等式结合，给出一个梯度的反向 Hölder 型估计。例如 \(d>2\) 时取 \(q=2d/(d+2)<2\)，有
\[
 \left(\frac1{|B_r|}\int_{B_r}|\nabla u|^2\right)^{1/2}
 \le C\left(\frac1{|B_{2r}|}\int_{B_{2r}}|\nabla u|^q\right)^{1/q}.
\]
这里积分除以球的体积，均取平均值。证明中先用能量估计将左端界为
\[
 \frac Cr\left(\frac1{|B_{2r}|}\int_{B_{2r}}|u-u_{B_{2r}}|^2\right)^{1/2},
\]
再用 Sobolev--Poincaré 将这一项界为右端。二维可选任意固定 \(1<q<2\)，由对应有限指数嵌入及 Hölder 不等式得到同样形式。

\ProseParagraphBreak

通常 Hölder 不等式以较高指数范数控制较低指数范数；这里对解的梯度出现相反方向，代价是扩大球。它暗示梯度具有比任意 \(L^2\) 函数更强的分布控制。由此继续推出高于 \(2\) 的可积指数，需要反向 Hölder 不等式的自改进定理；本节没有证明这一步，因而也不把更高可积性写成已经建立的结论。

\ProseParagraphBreak

局部有界性的迭代使用另一种增强：把能量估计应用于超过某个高度的部分。称 \(u\in H^1\) 为弱次解，是指 \(\int A\nabla u\cdot\nabla\varphi\le0\) 对所有非负紧支集测试函数成立。取 \(w=(u-k)_+\)，经截断逼近可用 \(\eta^2w\) 测试，得到
\[\int\eta^2|\nabla w|^2\le C\int w^2|\nabla\eta|^2.\]
在 \(\{u\le k\}\) 上 \(\nabla w=0\)，在其余部分 \(\nabla w=\nabla u\)，所以与前面的吸收步骤完全相同。高度 \(k\) 可以随迭代改变，这为同时控制高值集合的大小和函数值提供了自由度。

\subsection{局部有界性的动机}
\label{b2:subsec:160403}

\begin{theorem}[弱次解的局部上界]
\label{b2:thm:16-local-boundedness}
设 \(d\ge2\)，\(A\) 在 \(B_R\) 上有界可测且一致椭圆，\(u\in H^1(B_R)\) 是实弱次解。则
\[\operatorname*{ess\,sup}_{B_{R/2}}u_+
 \le C\left(\frac1{|B_R|}\int_{B_R}u_+^2\right)^{1/2},\]
其中 \(C\) 只依赖 \(d,\lambda,\Lambda\)。齐次弱解同时满足相应的 \(|u|\) 局部上界。
\end{theorem}
\begin{proof}
先缩放到 \(R=1\)。对支撑在单位球的 \(H_0^1\) 函数，选 \(\chi=d/(d-2)>1\)（\(d>2\)）或 \(\chi=2\)（\(d=2\)），Sobolev 不等式给出 \(\lVert v\rVert_{2\chi}^2\le C\lVert\nabla v\rVert_2^2\)。固定待选高度 \(K>0\)，令
\[r_n=\tfrac12+2^{-n-1},\qquad k_n=K(1-2^{-n}),\qquad
 Y_n=\int_{B_{r_n}}(u-k_n)_+^2.\]
选截断 \(\eta_n=1\) 于 \(B_{r_{n+1}}\)，支集在 \(B_{r_n}\)，且 \(|\nabla\eta_n|\le C2^n\)。令 \(w_n=(u-k_n)_+\)，截断能量估计连同乘积导数公式给出
\[\lVert\eta_nw_n\rVert_{2\chi}^2
 \le C\lVert\nabla(\eta_nw_n)\rVert_2^2\le C4^nY_n.\]
在集合 \(E_{n+1}=B_{r_{n+1}}\cap\{u>k_{n+1}\}\) 上，\(w_n\ge k_{n+1}-k_n=K2^{-n-1}\)。所以 \(|E_{n+1}|\le4^{n+1}K^{-2}Y_n\)。再用 Hölder 不等式，记 \(\theta=1-1/\chi>0\)，得
\[Y_{n+1}\le\lVert\eta_nw_n\rVert_{2\chi}^2|E_{n+1}|^\theta
 \le Cb^nK^{-2\theta}Y_n^{1+\theta},\]
其中 \(b>1\) 为固定常数。

令 \(Z_n=Y_n/K^2\)，则 \(Z_{n+1}\le Cb^nZ_n^{1+\theta}\)。若
\[Z_0\le C^{-1/\theta}b^{-1/\theta^2},\]
归纳可证 \(Z_n\le C^{-1/\theta}b^{-n/\theta-1/\theta^2}\)：代入递推式后，常数幂及 \(b\) 的指数恰好得到下一项。因此 \(Y_n\to0\)。取 \(K\) 为 \(\lVert u_+\rVert_{L^2(B_1)}\) 的足够大固定倍数即可满足初始小量条件；零范数情形结论显然。因为 \((u-K)_+^2\) 在 \(B_{1/2}\) 上的积分不超过每个 \(Y_n\)，故 \(u\le K\) 几乎处处。缩放回半径 \(R\) 得到平均形式的估计。对于解，\(u\) 与 \(-u\) 都是次解，分别估计正部即可。
\end{proof}

这一证明称为截断迭代。它没有控制系数的导数，也没有先假设解有界；有界性由高值集合的递推衰减产生。迭代中的超线性指数 \(1+\theta>1\) 来自 Sobolev 积分增益，正是小量逐步改善的原因。若只得到线性递推 \(Y_{n+1}\le Cb^nY_n\)，几何增长的常数会使论证失败。

\ProseParagraphBreak

一维可以直接使用 \(H^1\) 函数的绝对连续代表和基本定理获得局部控制，无须上述多维迭代。对于含源项的方程，截断估计会多出源项积分；要保留局部上界，源项需具有与维数相容的可积性。例如一般高维 \(L^2\) 源项未必足够。这里的齐次假设不能只从定理标题中读过而在应用时忘掉。

\subsection{De Giorgi–Nash–Moser 理论简介}
\label{b2:subsec:160404}

局部有界性还没有保证振荡随半径缩小而消失。实标量一致椭圆方程的更强结论是：非负解在较小球内的最大值与最小值可相互比较，任意实解因而具有 Hölder 连续代表。De Giorgi、Nash 和 Moser 的工作从不同方向建立了这一低正则理论。本小节介绍结论及其逻辑联系；其中 Harnack 比较的深层步骤不冒充为上一迭代已经证明。

\begin{theorem}[可测系数的 Harnack 比较，深层结论]
\label{b2:thm:16-measurable-harnack}
设 \(A\) 在 \(B_{2R}\subset\mathbb R^d\) 上为实有界可测矩阵，满足固定的一致椭圆上下界。若 \(u\in H^1(B_{2R})\) 为非负弱解 \(-\operatorname{div}(A\nabla u)=0\)，则
\[\operatorname*{ess\,sup}_{B_R}u\le C\operatorname*{ess\,inf}_{B_R}u,\]
其中 \(C\) 只依赖维数和椭圆上下界。
\end{theorem}
\begin{proofsketch}
上一小节已经完整证明次解的局部上界。余下的核心是超解的正性传播：一定比例的点上具有正下界，可以转成较小区域上的统一正下界。用负幂或对数测试建立这一性质，再结合覆盖分解，可以得到一个小正指数的弱 Harnack 下界；将它与相应指数的局部上界结合，就得到所列比较。这里省略正性传播和覆盖迭代的证明，故这一段只是证明路线。Dongsheng Li 与 Kai Zhang 的论文 \cite[Theorem 1.1 and Section 2]{Li2017Harnack} 给出可测系数的比较定理，并说明局部上界、正性传播与比较之间的关系；系统理论可参阅 \cite[Chapter 8]{Gilbarg2001Elliptic}。本章此前的 \(H^2\) 定理和局部有界性均不依赖此深层结论。
\end{proofsketch}

为了看出比较如何控制振荡，设 \(u\) 是 \(B_{2R}\) 内的任意有界实解，记在外球上的本性上、下界为 \(M,m\)，在内球上的为 \(M_1,m_1\)。函数 \(u-m\) 与 \(M-u\) 均为非负解。对二者应用 Harnack 比较，并将两项不等式相加，得到
\[(M-m)+(M_1-m_1)\le C\bigl((M-m)-(M_1-m_1)\bigr).\]
因此
\[\operatorname*{ess\,osc}_{B_R}u\le\vartheta\operatorname*{ess\,osc}_{B_{2R}}u,
 \qquad\vartheta=\frac{C-1}{C+1}<1.\]
在同心球上迭代，得到振荡按半径的某个正幂衰减。取每点缩小球平均值的极限，振荡衰减保证极限存在；重叠球比较保证这些值满足局部 Hölder 估计，并与原函数几乎处处一致。若 \(C=1\)，解在内球为常数，结论更直接。

\ProseParagraphBreak

这套结论的对象是实标量二阶散度方程。一般椭圆系统缺少同样的序结构，不能直接对向量解作正部截断，也不能把标量 Harnack 定理逐分量套用到耦合系统上。系数虽可跳跃，标量解仍可能 Hölder 连续；其梯度却未必连续，甚至不一定有平方可积的二阶导数。低正则理论改善的是解的振荡，与前面依赖系数导数的高阶 Sobolev 理论相互补充。

% 本章小结（不计入编号）
\begin{chaptersummary}
内部高阶正则性从一个与步长无关的差商估计开始。系数的 Lipschitz 控制保证通量差商的误差有界，而更高系数导数保证自举时产生的右端仍在所需空间。Sobolev 嵌入再把积分意义的导数转成经典代表，其指数门槛依赖维数。边界估计还消耗坐标图的正则性；凹角例子说明内部光滑不自动带来全局 \(H^2\)。

\ProseParagraphBreak

仅有界可测系数时，Caccioppoli 不等式仍然成立。截断迭代把高值集合的能量与测度联系起来，完整给出齐次标量弱解的局部有界性。Harnack 比较进一步控制振荡，但需要另行建立超解的正性传播；从比较到 Hölder 衰减的推论则可直接迭代得到。至此，变分存在性、弱型可解性与正则性三项任务已各有明确的假设和证明。
\end{chaptersummary}

比较本编与后面的估计时，可以固定一个实标量 Dirichlet 问题，再逐项查看增加了什么条件、换来了什么结论。\cref{b2:tab:16-dirichlet-levels}中的各行保留自己的空间和假设，并不是把所有条件一并加到同一个定理上。

\begin{table}[htbp]
\centering
\small
\caption[椭圆 Dirichlet 问题的各层结论]{椭圆 Dirichlet 问题的表示、存在性与正则性}
\label{b2:tab:16-dirichlet-levels}
\begin{tabularx}{\textwidth}{@{}p{0.16\textwidth}XX@{}}
\toprule
所处阶段 & 已得到什么 & 应保留什么条件或限制 \\
\midrule
第十三章 & 球、半空间等模型的 Green 表示与 Poisson 公式。 & 表示一个已有解，与证明一般区域上存在解，是不同任务；核、数据及边界取值须满足对应公式的条件。 \\
第十四章 & 在给定仿射边界类中极小化凸能量。 & 二次能量使用 \(H_0^1\)、合法的边界延拓及 \(H^{-1}\) 源项；非二次能量另核增长、自反性及弱下半连续性。 \\
第十五章 & 强制双线性型的弱解存在、唯一及能量估计。 & 系数有界可测且一致椭圆，源项及测试处于相应对偶空间；不要求系数对称，但不能省去强制性或边界类。 \\
第十六章光滑系数部分 & Lipschitz 系数给内部 \(H^2\)；满足边界条件时进一步给全局 \(H^2\)。 & 内部只考察紧含子域；全局另用有界 \(C^{1,1}\) 区域、\(L^2\) 源项，非齐次边界须有 \(g\in H^2\) 且 \(\operatorname{Tr}u=\operatorname{Tr}g\)。凹角会破坏全局结论。 \\
第十六章及附录 H.1 & 对相应齐次实标量方程作 Harnack 比较并得局部 Hölder 性。 & 比较用于非负弱解和有界可测的一致椭圆系数；不依赖系数可微性，也不给出同样的二阶导数结论。 \\
第十八、二十二章 & 全空间常系数二阶导数的 \(L^p\) 控制，以及局部化、冻结系数的步骤。 & \(1<p<\infty\)；齐次导数估计不自动控制整个非齐次 \(W^{2,p}\) 范数。边界、变系数误差及自举起点仍须另核。 \\
\bottomrule
\end{tabularx}
\end{table}

同一个 \(H^1\) 弱解可以在内部光滑，却因粗糙边界数据或区域凹角而不在全局 \(H^2\)。反过来，对只是可测的系数，Harnack 理论给出的正性与振荡控制，也不是差商高阶理论的简写。\secref{b2:sec:1805}与\secref{b2:sec:2203}将把二阶导数估计放到奇异积分、乘子的框架中；它们增加的是可积指数的分析工具，不撤销已有的边界与系数条件。


\BookExercises

各题均为本书自拟。题目分别改变尺度、系数、数据与边界，检验不同正则性结论的边界；有界解的 Liouville 选做题使用本章明示、附录 H 证明的 Harnack 比较。连续边值题综合变分、弱形式与边界数据，说明内部光滑不等于全局二阶正则；最后的选做题训练边界坐标与通量恢复的具体计算。

% 自拟：差商平移约定与离散能量恒等式。
\Needspace{6\baselineskip}
\begin{exercise}\label{b2:ex:16-discrete-energy}
设 \(d\ge1\)、\(1\le k\le d\)，\(u\in H^1(\mathbb R^d;\mathbb R)\)，\(D_h^ku(x)=[u(x+he_k)-u(x)]/h\)。证明当非零实数 \(h\to0\) 时 \(D_h^ku\to\partial_ku\) 于 \(L^2\)，并证明
\[\int_{\mathbb R^d}u(-D_{-h}^kD_h^ku)=\lVert D_h^ku\rVert_2^2.\]
说明为何差商测试中使用负的 \(D_{-h}^k\)。
\end{exercise}
\begin{solution}
对光滑紧支集函数，线段积分公式和导数的平移连续性给出强收敛。一般 \(u\) 先取 \(H^1\) 光滑逼近，利用差商统一界 \(\lVert D_h^k(u-v)\rVert_2\le\lVert\partial_k(u-v)\rVert_2\)，三角不等式把结论推广。恒等式由离散分部积分得到：\(\int u(-D_{-h}^kw)=\int(D_h^ku)w\)，取 \(w=D_h^ku\) 即可。若漏掉负号，主项就成为负的平方范数，无法使用椭圆正性吸收误差。
\end{solution}

% 自拟：系数Lipschitz但不具二阶正则性的明确反例。
\Needspace{6\baselineskip}
\begin{exercise}\label{b2:ex:16-lipschitz-not-h3}
在 \((-1,1)\) 上令 \(u(x)=\int_0^x(1+|t|)^{-1}\,\mathrm dt\)。证明 \(-((1+|x|)u')'=0\) 弱成立，\(u\in H^2(-1,1)\)，但 \(u\notin H^3(-1,1)\)。
\end{exercise}
\begin{solution}
通量恒为 \(1\)，故与任意测试函数导数的积分为零。\(u'=(1+|x|)^{-1}\) 连续且 Lipschitz，其弱导数在两侧分别为 \((1-x)^{-2}\)（\(x<0\)）与 \(-(1+x)^{-2}\)（\(x>0\)），有界，故 \(u\in H^2\)。\(u''\) 从 \(1\) 跳到 \(-1\)，第三分布导数包含 \(-2\delta_0\)，不能属于 \(L^2\)。这说明右端光滑不足以越过系数正则性的限制。
\end{solution}

% 自拟：全空间常系数二阶估计的Fourier核验。
\Needspace{6\baselineskip}
\begin{exercise}\label{b2:ex:16-fourier-h2}
设 \(f\in L^2(\mathbb R^d)\)。用本卷的 \(2\pi\) Fourier 约定求 \((I-\Delta)u=f\) 的唯一 \(L^2\) 分布解，并证明 \(u\in H^2\)、\(\lVert u\rVert_{H^2}\le C_d\lVert f\rVert_2\)。
\end{exercise}
\begin{solution}
定义 \(\widehat u(\xi)=\widehat f(\xi)/(1+4\pi^2|\xi|^2)\)。该函数属于 \(L^2\)，Plancherel 给出 \(u\)，且乘子恒等式验证分布方程。对 \(|\alpha|\le2\)，函数 \((2\pi i\xi)^\alpha/(1+4\pi^2|\xi|^2)\) 有界，故弱导数 \(D^\alpha u\in L^2\)，并有相应范数界；对有限个指标求和得到估计。若另一 \(L^2\) 解存在，差的 Fourier 变换被严格正的 \(1+4\pi^2|\xi|^2\) 乘后为零，故差为零。
\end{solution}

% 自拟：L2右端在高维不足以保证局部有界。
\Needspace{6\baselineskip}
\begin{exercise}\label{b2:ex:16-l2-unbounded}
在 \(\mathbb R^5\) 中取 \(\chi\in C_c^\infty(B_1)\)，在原点附近为 \(1\)，令 \(u(x)=\chi(x)|x|^{-1/4}\)。证明 \(u\in H_0^1(B_1)\)、\(-\Delta u\in L^2(B_1)\)，但 \(u\) 无界。解释为何这不违反本章局部有界性定理。
\end{exercise}
\begin{solution}
原点附近 \(|u|\) 量级为 \(r^{-1/4}\)，梯度量级为 \(r^{-5/4}\)；乘上五维体积因子 \(r^4\) 后，两者平方均可积。截去一个小球并平滑，误差的 \(H^1\) 范数趋零，故紧支集函数 \(u\) 属于 \(H_0^1\)。由 \(\Delta r^\alpha=\alpha(\alpha+3)r^{\alpha-2}\)，主奇性为 \(r^{-9/4}\)，其平方积分近零为 \(\int r^{-1/2}\,\mathrm dr<\infty\)。小球边界通量及测试函数误差趋零，故无额外点质量。于是 \(-\Delta u\in L^2\)，但 \(u\to\infty\)。局部有界性定理讨论齐次方程，此例右端非零；高维中 \(L^2\) 源项不足以支持同一结论。
\end{solution}

% 自拟：角点指数与Sobolev阈值。
\Needspace{6\baselineskip}
\begin{exercise}\label{b2:ex:16-corner-threshold}
设 \(0<\omega<2\pi\)，\(\alpha=\pi/\omega\)，\(v=r^\alpha\sin(\alpha\theta)\) 定义在平面扇形。对 \(\alpha\ne1\)，求它在角点附近属于 \(W^{2,p}\)、\(1\le p<\infty\) 的准确条件，并单独解释 \(\alpha=1\) 的情形。
\end{exercise}
\begin{solution}
若 \(\alpha\ne1\)，二阶导数的平方和在角向上有非零有界因子，径向量级为 \(r^{\alpha-2}\)。因此二阶导数的 \(p\) 次积分有限恰当且仅当 \(p(\alpha-2)+1>-1\)，即 \(p(2-\alpha)<2\)。此条件也保证较低阶项可积；等号时对数发散。若 \(\alpha\ge2\)，该严格不等式对一切有限 \(p\) 自动成立。\(\alpha=1\) 时 \(v=r\sin\theta=y\) 是线性函数，二阶导数恒零，不能用只看形式幂次的阈值排除它。
\end{solution}

% 自拟：迭代引理的显式小量阈值。
\Needspace{6\baselineskip}
\begin{exercise}\label{b2:ex:16-iteration-lemma}
设非负实数列 \((Y_n)_{n\ge0}\) 满足 \(Y_{n+1}\le Ab^nY_n^{1+\delta}\)，其中 \(A>0\)、\(b>1\)、\(\delta>0\)。证明当 \(Y_0\le A^{-1/\delta}b^{-1/\delta^2}\) 时，\(Y_n\le A^{-1/\delta}b^{-n/\delta-1/\delta^2}\)。说明仅知道 \(Y_0<1\) 为何不总够用。
\end{exercise}
\begin{solution}
初始步正是给定条件。若第 \(n\) 步成立，代入递推式得到 \(A\) 的幂为 \(1-(1+\delta)/\delta=-1/\delta\)，\(b\) 的幂为 \(n-(1+\delta)n/\delta-(1+\delta)/\delta^2=-(n+1)/\delta-1/\delta^2\)，故归纳闭合。特别地序列趋零。小量阈值依赖 \(A,b,\delta\)；若 \(A\) 很大，某个接近 \(1\) 的初值即使小于 \(1\)，第一次递推也可能增大，所以必须通过选择截断高度或归一化满足定量阈值。
\end{solution}

% 自拟：全空间平方可积齐次解的消失。
\Needspace{6\baselineskip}
\begin{exercise}\label{b2:ex:16-l2-liouville}
设 \(d\ge2\)，\(A\in L^\infty(\mathbb R^d;\mathbb R^{d\times d})\)，对几乎处处 \(x\) 和全部 \(\xi\in\mathbb R^d\) 有 \(|A(x)\xi|\le\Lambda|\xi|\)、\(\xi^{\mathsf T}A(x)\xi\ge\lambda|\xi|^2\)，其中 \(0<\lambda\le\Lambda<\infty\) 不随 \(x\) 变化。若实函数 \(u\in H_{\mathrm{loc}}^1(\mathbb R^d)\cap L^2(\mathbb R^d)\) 弱解 \(-\operatorname{div}(A\nabla u)=0\)，不用系数可微性，证明 \(u=0\) 几乎处处。
\end{exercise}
\begin{solution}
对任意中心和半径 \(R\)，局部有界性给出 \(\operatorname*{ess\,sup}_{B_{R/2}}|u|\le C_dR^{-d/2}\lVert u\rVert_{L^2(\mathbb R^d)}\)。固定任意有限球并使 \(R\to\infty\)，右端趋零，因此该有限球内 \(u=0\) 几乎处处。用可数个球覆盖全空间即可。这里全局共同的椭圆常数保证估计中的 \(C_d\) 不随半径增长。
\end{solution}

% 自拟：深层Harnack输入推出有界解的Liouville性质。
\Needspace{6\baselineskip}
\begin{optionalexercise}\label{b2:ex:16-bounded-liouville}
设 \(d\ge2\)，\(A\in L^\infty(\mathbb R^d;\mathbb R^{d\times d})\)，对几乎处处 \(x\) 和全部 \(\xi\in\mathbb R^d\) 有 \(|A(x)\xi|\le\Lambda|\xi|\)、\(\xi^{\mathsf T}A(x)\xi\ge\lambda|\xi|^2\)，其中 \(0<\lambda\le\Lambda<\infty\) 为固定常数。若实函数 \(u\in H_{\mathrm{loc}}^1(\mathbb R^d)\cap L^\infty(\mathbb R^d)\) 弱解 \(-\operatorname{div}(A\nabla u)=0\)，使用本章的深层 Harnack 结论证明 \(u\) 为常数。明确指出本题比平方可积解的 Liouville 题多用了哪项输入。
\end{optionalexercise}
\begin{solution}
由正文的振荡收缩，存在与球无关的 \(0<\vartheta<1\)，使 \(\operatorname*{ess\,osc}_{B_R}u\le\vartheta^n\operatorname*{ess\,osc}_{B_{2^nR}}u\le2\vartheta^n\lVert u\rVert_\infty\)。令 \(n\to\infty\)，每个球内本性振荡为零。重叠球的常数相同，故全空间几乎处处为同一常数；取 Hölder 代表即处处为常数。本题使用 Harnack 导出的振荡收缩，上一题只用已经完整证明的局部上界与全局 \(L^2\) 衰减。
\end{solution}

% 自拟：综合圆周正交展开、凸 Dirichlet 能量、弱形式及内部/边界正则性的区别；光滑域上的障碍来自边界数据。
\Needspace{6\baselineskip}
\begin{exercise}[连续边值与二阶正则性的限制]\label{b2:ex:16-rough-trace-energy}
在单位圆盘 \(D=\{x\in\mathbb R^2:|x|<1\}\) 上取实函数空间。令
\[
 g(\theta)=\sum_{n=1}^\infty n^{-3/2}\cos(n\theta),\qquad
 u(r,\theta)=\sum_{n=1}^\infty n^{-3/2}r^n\cos(n\theta),\quad 0\leq r<1.
\]
在原点按上述各调和多项式的值定义 \(u(0)=0\)。
\begin{enumerate}
\item\label{b2:item:16-trace-energy-class}证明 \(g\) 连续、\(u\in H^1(D)\cap C(\overline D)\)，且 \(u\) 的 Sobolev 迹等于 \(g\)。计算 \(\int_D|\nabla u|^2\)。
\item\label{b2:item:16-trace-energy-minimum}证明 \(u\) 是凸能量 \(E(v)=\frac12\int_D|\nabla v|^2\) 在仿射类 \(u+H_0^1(D)\) 中的唯一极小元，并写出对全部 \(H_0^1(D)\) 测试函数成立的弱方程。
\item\label{b2:item:16-trace-energy-regularity}证明 \(u\) 在圆盘内部光滑，但 \(u\notin H^2(D)\)。解释这为何不与光滑区域上的全局 Dirichlet 二阶正则性矛盾。
\end{enumerate}
\end{exercise}
\begin{solution}
因为 \(\sum_n n^{-3/2}<\infty\)，两级数分别在圆周和闭圆盘上一致绝对收敛；\(r^n\cos(n\theta)\) 是 \(\operatorname{Re}(x_1+\mathrm ix_2)^n\)，故不存在原点的坐标奇性。于是 \(u\) 连续到边界，边界值为 \(g\)。记其前 \(N\) 项为 \(u_N\)。圆周正交性和极坐标积分给出
\[
 \|u_N\|_{L^2(D)}^2
 =\frac\pi2\sum_{n=1}^N\frac{n^{-3}}{n+1},\qquad
 \|\nabla u_N\|_{L^2(D)}^2
 =\pi\sum_{n=1}^N n^{-2}.
\]
将相同计算用于两个部分和的差，得到 \(H^1\) Cauchy 性。其 \(L^2\) 极限与一致极限相同，所以 \(u\in H^1(D)\)，且
\(\int_D|\nabla u|^2=\pi\sum_{n=1}^\infty n^{-2}\)。圆盘是光滑有界区域，\subsecref{b2:subsec:150103}所用的迹算子在 \(H^1\to L^2(\partial D)\) 间连续；\(u_N\to u\) 于 \(H^1\)，而边界部分和一致趋于 \(g\)，因此 \(\operatorname{Tr}u=g\)。该小节的零迹刻画还说明，\(u+H_0^1(D)\) 恰是全部具有这一 Sobolev 迹的 \(H^1\) 函数。

每个 \(u_N\) 都调和，因而对 \(\varphi\in C_c^\infty(D)\)，有 \(\int_D\nabla u_N\cdot\nabla\varphi=0\)。由 \(H^1\) 收敛取极限，再由 \(C_c^\infty(D)\) 在 \(H_0^1(D)\) 中的稠密性，得到
\[
 \int_D\nabla u\cdot\nabla w=0\qquad(w\in H_0^1(D)).
\]
所以对每个允许扰动 \(w\)，
\[
 E(u+w)=E(u)+\int_D\nabla u\cdot\nabla w
                 +\frac12\|\nabla w\|_2^2
 =E(u)+\frac12\|\nabla w\|_2^2.
\]
这既证明极小性，也由 \(H_0^1\) 上的 Poincaré 不等式证明唯一性：取等要求 \(\nabla w=0\)，从而 \(w=0\)。弱方程、固定迹的变分问题和凸性在这里分别负责识别候选者、限定允许扰动与判断唯一极小。

在任意 \(r\leq R<1\) 上，固定阶导数的级数由 \(\sum_n n^M R^{n-M}\) 的尾部控制，故 \(u\) 内部光滑且逐点调和。这也符合\cref{b2:cor:16-smooth-interior}：单位系数和零右端都光滑，且 \(u\) 已是能量弱解。可是对 \(0<R<1\)，仍由正交性，
\[
 \int_{|x|<R}|\partial_{rr}u|^2\,\mathrm dx
 =\frac\pi2\sum_{n=2}^\infty\frac{n-1}{n}R^{2n-2}.
\]
令 \(R\uparrow1\)，非负级数的极限为无穷。若 \(u\in H^2(D)\)，则几乎处处 \(x\ne0\) 有 \(\partial_{rr}u=e_r^{\mathsf T}D^2u\,e_r\)，其中 \(e_r=x/|x|\)，其平方积分应受 Hessian 的平方积分控制，矛盾。圆盘本身光滑，失败来自边界数据。\cref{b2:thm:16-boundary-h2}要求零 Dirichlet 迹，而此处 \(g\ne0\)；非齐次版本\cref{b2:cor:16-nonhomogeneous-boundary-h2}另要求有 \(H^2\) 的边界提升，不能仅凭边值连续就套用。事实上，若这样的提升存在，该推论就会迫使 \(u\in H^2(D)\)，与上述计算矛盾。
\end{solution}

% 自拟：二维图像坐标下计算几何系数，以通量恢复纯法向导数；不承担正文证明依赖。
\Needspace{8\baselineskip}
\begin{optionalexercise}[边界平坦化与法向通量]\label{b2:ex:16-flattening-flux}
设实函数 \(\psi\in C^{1,1}([-2,2])\)，记 \(M=\|\psi'\|_\infty\)，并令
\[
 Q=(-2,2)\times(0,2),\quad Q_1=(-1,1)\times(0,1),\qquad
 F(s,t)=(s,t+\psi(s)),\quad \Omega_\psi=F(Q).
\]
设 \(f\in L^2(\Omega_\psi;\mathbb R)\)，\(u\in H^1(\Omega_\psi;\mathbb R)\) 弱解 \(-\Delta u=f\)，且在下边界图像 \(F((-2,2)\times\{0\})\) 取零迹。置 \(v=u\circ F\)、\(\widetilde f=f\circ F\)。
\begin{enumerate}
\item\label{b2:item:16-flux-coordinate}计算 \(DF\)、\(\det DF\)，从弱形式导出变换系数
\[
 \widetilde A=\begin{pmatrix}1&-\psi'\\-\psi'&1+(\psi')^2\end{pmatrix}.
\]
证明 \((2+M^2)^{-1}|\xi|^2\le\xi^{\mathsf T}\widetilde A\xi\le(2+M^2)|\xi|^2\)，并说明变换后的测试函数在平坦边界及其余边上的要求。
\item\label{b2:item:16-flux-recovery}假设切向差商已给出 \(v_{ss},v_{st}\in L^2(Q_1)\)。不预设 \(v_{tt}\) 存在，使用 \(q=(1+(\psi')^2)v_t\) 先证明 \(\partial_tq\in L^2(Q_1)\)，再求出 \(v_{tt}\) 的公式和 \(L^2\) 上界。
\item\label{b2:item:16-flux-geometry}在以上论证中，\(\psi'\) 有界和 \(\psi''\) 本性有界各用在何处？为何不能只凭边界为 Lipschitz 就沿用同一二阶估计？
\end{enumerate}
\end{optionalexercise}
\begin{solution}
有
\[
 DF=\begin{pmatrix}1&0\\\psi'&1\end{pmatrix},\qquad
 (DF)^{-1}=\begin{pmatrix}1&0\\-\psi'&1\end{pmatrix},\qquad \det DF=1.
\]
对 \(\zeta\in C_c^\infty(Q)\)，令 \(\varphi=\zeta\circ F^{-1}\)。它属于 \(H_0^1(\Omega_\psi)\)，可由密度用于原弱形式。链式法则和换元给出
\[
 \int_Q\widetilde A\nabla v\cdot\nabla\zeta
 =\int_Q\widetilde f\,\zeta,\qquad
 \widetilde A=(DF)^{-1}(DF)^{-\mathsf T}.
\]
计算矩阵乘积便得题中系数。该实对称矩阵行列式为 \(1\)，迹为 \(2+(\psi')^2\)，且二次型为 \((\xi_s-\psi'\xi_t)^2+\xi_t^2\)，故两个特征值为正，上界不超过 \(2+M^2\)，下界由两者乘积为 \(1\) 得到。变换后的 \(v\) 在 \(t=0\) 上取零迹。切向差商所用测试函数可以接触此平坦边界，但必须在那里取零迹，并在左右边和上边附近消失；这样其逆变换仍在原区域的 \(H_0^1\) 中。

记 \(p(s)=\psi'(s)\)。方程的两个通量分量为 \(v_s-pv_t\) 和 \(-pv_s+(1+p^2)v_t\)。已有混合导数的分布交换律给 \(\partial_tv_s=v_{st}\)，所以
\[
 \partial_tq
 =-\widetilde f-\partial_s(v_s-pv_t)+\partial_t(pv_s)
 =-\widetilde f-v_{ss}+p'v_t+2pv_{st}\in L^2(Q_1).
\]
这里只对已有相应方向 \(L^2\) 弱导数的因子使用该方向的乘积法则：\(\partial_sv_t=v_{st}\)、\(\partial_tv_s=v_{st}\) 已知，而 \(v_{tt}\) 尚未得到。因为 \((1+p^2)^{-1}\) 是有界的 \(W^{1,\infty}\) 函数且不依赖 \(t\)，由 \(v_t=(1+p^2)^{-1}q\) 的法向弱微分得到
\[
 v_{tt}=\frac{-\widetilde f-v_{ss}+\psi''v_t+2\psi'v_{st}}{1+(\psi')^2},
\]
以及
\[
 \|v_{tt}\|_{L^2(Q_1)}\le
 \|\widetilde f\|_{L^2(Q_1)}+\|v_{ss}\|_{L^2(Q_1)}
 +\|\psi''\|_\infty\cdot\|v_t\|_{L^2(Q_1)}+2M\|v_{st}\|_{L^2(Q_1)}.
\]
\(\psi'\) 的上界控制坐标的一阶范数和椭圆常数；\(\psi''\) 的上界控制几何系数的一阶导数，也控制最后的 \(\psi''v_t\)。只有 Lipschitz 图像时，\(\psi'\) 虽本性有界，其弱导数未必为有界函数，切向系数差商和上述乘积界便缺少依据。因此斜率受控与斜率变化受控在二阶正则性中承担不同职责。
\end{solution}

